\documentclass[../main.tex]{subfiles}
\begin{document}
%==========================================TIÊU ĐỀ TẬP========================================
\begin{table}[h]
    \begin{adjustwidth}{-1cm}{}
        \begin{tabular}{>{\centering\arraybackslash}p{6cm}|>{\raggedright\arraybackslash}p{14cm}}
            \multirow{5}{6cm}
            % Cột bên trái: ÉPISODE và số tập
            {\\[-40pt]
                \flaregothic\fontsize{20pt}{20pt}\selectfont \centering
                \color{eptype}HISTOIRE\color{black}\\
                \gangofthree\fontsize{72pt}{72pt}\selectfont
                \color{epnum}4
                \\[18pt]
                % \flaregothic\fontsize{14pt}{14pt}\selectfont
                % \color{epattr}BIENVENUE, \uline{\color{Banpire} Mel} !
                % \flaregothic\fontsize{18pt}{18pt}\selectfont
                % \color{epattr}ÉPISODE PERDU
                \color{black}
            }
             &
            % Dòng thứ nhất: tên tập tiếng Nhật
            {\fotskip\fontsize{18pt}{18pt}\selectfont \ruby{波}{は}\ruby{乱}{らん}の\ruby{大}{おお}\ruby{晦}{みそ}\ruby{日}{か}！}
            \\[6pt]
            % Dòng thứ hai: tên tập tiếng Anh
             & {\cafeteria\fontsize{18pt}{18pt}\selectfont A Chaotic New Year's Eve!}      \\[4pt]
            % Dòng thứ ba: tên tập tiếng Việt
             & {\vnmsans\fontsize{18pt}{18pt}\selectfont Ngày cuối năm đầy bất ổn!}        \\[8pt]
            % Dòng thứ tư: Nhân vật xuất hiện
             & {\sen{\fontsize{14pt}{14pt}\selectfont Track used: Firecrackers 2 - F-777}}
        \end{tabular}
    \end{adjustwidth}
\end{table}
%===========================================NỘI DUNG==========================================
\color{black}

\pov{Hikawa Kuon}

Sáng ngày 30 Tết, còn 19 tiếng nữa là đến đêm giao thừa.

Như thường lệ, tôi giữ thói quen dậy sớm vào lúc năm giờ sáng, thói quen được hình thành
từ những năm cấp hai và vẫn duy trì cho đến bây giờ, kể cả khi đó là ngày nghỉ.

Nếu các bạn nghĩ rằng hôm nay sẽ là một ngày nghỉ như bình thường, thì tôi xin đính
chính: hôm nay không hề bình thường, thậm chí còn chẳng phải là một ``ngày nghỉ'' đúng
nghĩa.

Sự khác biệt đã bắt đầu ngay từ lúc tôi thức dậy. Thông thường, chỉ có hai bố con tôi nằm
ngủ trên tấm đệm quen thuộc ở tầng năm, còn mẹ và em gái tôi thì ngủ ở tầng ba. Nhưng
sáng nay, tấm đệm này lại chật đến khó thở. Không chỉ có tôi, mà còn thêm\dots\ đúng vậy,
bốn thành viên từ nhóm \textit{GAMERS} của \hll.

Bạn thắc mắc tại sao họ lại ở đây? Tại sao họ lại ngủ chung với tôi? Chuyện là thế này\dots

\begin{center}
    \uline{Tối hôm qua.}
\end{center}

Sau khi cả nhà tôi và các cô gái cùng nhau dọn dẹp xong xuôi, mọi người quay về nơi trú
chân để nghỉ ngơi. Nói là ``nhà'' nhưng thật ra gia đình tôi đã ở sẵn đây, còn các cô gái và
A-san thì thuê căn nhà bên cạnh, vốn được cải tạo từ một quán cà phê.

Tôi không rõ họ ăn tối ở đâu, nhưng nghe nói là họ đi ăn ngoài. Có thể họ được đồng
nghiệp hoặc ông sếp Tanigo-san tài trợ chi phí, hoặc họ tự bỏ tiền túi. Cá nhân tôi nghiêng
về phương án thứ hai, vì chi phí ở đây rẻ hơn nhiều so với Tokyo hoặc thành phố Điện tử,
ít nhất cũng phải ba lần.

Trong khi đó, gia đình tôi chỉ ăn một bữa cơm tối đơn giản, giống như mọi ngày. Không
quá thanh đạm nhưng cũng chẳng cầu kỳ, vừa đủ để tận hưởng không khí ấm cúng. Với
tôi, cơm nhà nấu vẫn luôn là ngon nhất.

Mọi việc sau đó diễn ra khá bình thường, giống như mọi ngày. Bố tôi nằm vừa đọc sách
vừa nghe tivi, tôi và em gái chơi điện tử, còn mẹ tôi thì bận rộn với vài công việc vặt trong
nhà.

Nhưng rồi bất chợt, con Mi nhà tôi chạy ra cửa, đuôi vẫy lia lịa và sủa oăng oẳng không
ngừng.

``{\color{red} Gì thế, Mi? Có cái gì mà mày sủa oăng oẳng thế?}'' tôi vừa nói vừa bước ra, nhưng câu
nói còn chưa dứt thì đã thấy bốn bóng người từ từ bước lên cầu thang dẫn vào phòng của
tôi.

Ban đầu, tôi còn tưởng đó là nhân viên của cửa hàng viễn thông chạy lên để lấy thang.
Tuy nhiên, nhìn cách mà con chó nhà tôi sủa, tôi dám chắc rằng họ là những người lạ, ít
nhất là đối với nó.

Tôi còn nghe loáng thoáng tiếng bốn cô gái từ dưới nhà lên. Những giọng nói nghe rất
quen thuộc\dots

Một giọng nói tỏ vẻ nghiêm túc cất lên: ``Tôi không nghĩ đó là ý hay đâu\dots'' Trong khi đó,
một giọng nói khác tỏ vẻ vui tươi hơn cũng cất lên: ``Thôi nào Mio, tôi nghĩ Kuon-senpai
sẽ đông ý thôi\~{}''

Một giọng nói khác lanh lảnh vang lên, giống như một con chó biết nói sẽ reo lên: ``Kuon-senpai\~{}!
Anh có ở nhà không?'' Thì cũng phải thôi, vì đó là Korone-chan mà, và tôi biết
ngay mọi chuyện sẽ không hề ``bình thường'' một chút nào.

\dots Đúng, là họ đang lên. Nhóm \textit{GAMERS}, với nàng Sói đen, nàng Cáo trắng, và cặp đôi
Chó-Mèo. Bốn vị khách ``không mời mà đến'', với mỗi người mang sẵn cả bộ đồ ngủ.

Fubu-chan nhanh nhảu cất tiếng: ``Cháu xin phép ạ!'', trong khi Mio-chan lững thững theo
sau, lí nhí nói: ``C-Cháu xin phép\dots'' Hai người ở sau cùng cũng đồng thanh hô đầy phấn
khởi, bằng tiếng bản xứ: ``{\color{red} Chúng cháu chào các bác ạ!}''

Bố mẹ tôi đáp lại lời chào khá nồng hậu: ``Ừ, các cháu vào đi.'' Riêng tôi thì chỉ có thể
cười trừ vì đã đoán được họ định sẽ làm gì ở nhà tôi.

Mẹ tôi lên tiếng: ``Các cháu lên chơi hả?'' cả bốn cô gái đều đồng thanh một tiếng ``Vâng!''
rất to.

\dots Tôi biết, lý do mà họ lên đây không phải là lên đây, mà vì họ muốn ngủ với tôi. . Thực
ra, tôi không ngại việc họ ngủ chung với một cậu con trai như thôi, vì tôi thế nào cũng
được, nhưng tại sao họ lại muốn ngủ ở một nơi chật hẹp thế này?

``Chuyện gì thế, mấy đứa?'' tôi hỏi, tay xoa đầu con Mi đang vẫy đuôi mừng rỡ như thể
nó vừa gặp người quen cũ.

``Bọn em không ngủ được ở bên quán café,'' Okayu trả lời, gương mặt có vẻ khá bình
thản.

Tôi nghiêng đầu, tỏ vẻ hơi ngạc nhiên: ``Không ngủ được? Quán ấy chẳng phải yên tĩnh
lắm sao?''

Cáo trắng chen vào, vẻ hơi bối rối: ``Dạ\dots\ Lúc đầu chỗ đó cũng yên lặng thật, nhưng rồi
Matsuri-chan và Haato-chan bắt đầu chơi bắn nước, rồi kéo cả Subaru-chan vào\dots\ Cuối
cùng thì gần như cả tầng đều náo loạn. Bọn em đành bỏ qua bên này thôi!''

Được đà, Korone-chan còn gật đầu lia lịa, thêm thắt: ``Anh thấy đấy, tụi em là nhóm
GAMERS, cần không gian yên tĩnh để tái tạo năng lượng. Nhà anh chắc chắn sẽ hợp
hơn!''

Đến lượt bố tôi hỏi, bất ngờ khi nhìn thấy họ mang cả đồ ngủ mang theo: ``Chú cũng đoán
được bên đó loạn thế nào rồi. Thế, chắc các cháu cũng định ngủ lại đây luôn nhỉ?''

Koro-chan đáp một hồn nhiên: ``Đúng đó chú!''

``Khoan, khoan, khoan\dots'' Tôi cố nhấn mạnh, lý giải với bốn người bọn họ: ``Mấy đứa
có thể nhờ A-san sắp xếp phòng khác mà?''

Mio-chan lên tiếng, tỏ vẻ bất lực: ``Chị ấy bảo phòng nào cũng đầy mất rồi. Bọn em tính
là ra ngoài thuê phòng trọ một đêm, nhưng mà\dots''

``\dots bọn em nghĩ rằng, thôi thì lên chỗ Kuon-senpai ngủ cho tiện!'' Fubuki cười, nhún vai
đáp. Bạn của em ấy khẽ lắc đầu ngao ngán.

Okayu-chan lặng lẽ vỗ vai tôi, gương mặt đầy vẻ thấu hiểu: ``Anh yên tâm, tụi em không
phá đâu. Chỉ cần góc nhỏ trên đệm là đủ.'' Bạn thân của em ấy cũng tiếp lời: ``Bọn em
ngủ gọn lắm, không chiếm nhiều chỗ đâu\~{}''

Tôi chỉ biết thở dài. ``Nhà cô chật thế này, mấy đứa có ngủ được không vậy?'' mẹ tôi hỏi,
đúng ý mà tôi đang định hỏi. Tôi cũng không tin lắm việc cả bốn người họ sẽ ``ngủ mà
không chiếm nhiều chỗ.''

Đoán được ý định của tôi, Fubuki-chan nhanh nhảu nói: ``Không sau đâu, cô chú ạ! Bọn
cháu đã nhiều lần cùng lúc ngủ một giường rồi mà, nên cháu nghĩ nơi này vẫn còn rộng
lắm!''

``Cháu với Koro-san cũng đã từng ngủ với nhau rồi ạ\~{}'' bạn đồng hành của nàng Khuyển
nói.

Chỉ trừ Sói đen-chan thì vẫn giữ vẻ nghiêm túc: ``Nhưng đấy là khi mọi người sang chơi
với nhau! Còn đây là nhà riêng của anh ấy, ta không thể tùy tiện được!''

Tôi chỉ có thể cười xòa, cố gắng dàn hòa: ``Không sao đâu, mấy đứa cứ thoải mái, miễn
đừng có phá nhà anh giống như cái đám `PON' ban nãy là được.''

Nghe tôi nói vậy, Cáo trắng-chan quay sang bạn của em ấy, nháy mắt: ``Thấy chưa Mio?
Anh ấy bảo là chúng ta cứ xõa đi rồi\~{}'' Dù thế, cô bạn nghiêm túc của em ấy vẫn lắc đầu
ngán ngẩm, tỏ vẻ bất lực với một Fubu-chan thảnh thơi: ``Nhưng ở đây có mỗi một cái
đệm, vậy chúng ta sẽ ngủ thế nào?''

``Chuyện đó để chú lo,'' bố tôi đáp, ánh mắt rộ lên vẻ tinh quái. ``Mẹ nó xuống ngủ với
con Cún, còn chú sẽ qua phòng bác Hijaki ngủ với bà nội.''

``Con ngủ với bốn người con gái thế này, sướng quá còn gì nữa,'' mẹ tôi nói, buông giọng
trêu chọc.

Tôi chỉ biết lắc đầu với bố mẹ rất cơ hội đó. Chẳng biết tôi có thể ngủ ngon không, nhưng
mà một mình tôi với bốn người con gái\dots\ Có lẽ sẽ là một trải nghiệm gì đó khá kỳ cục.

Và thế là cả tối hôm đó, tôi dành thời gian với nhóm \textit{GAMERS}. Vì họ là \textit{GAMERS}, nên
hiển nhiên là họ sẽ chơi game rồi. Cậy nhà tôi có chiếc ti-vi, thế là cả bốn người chơi một
vài trò đồng đội trên S*itch. Hết thảy ba tiếng đồng hồ, từ tám giờ tối đến đêm. Cả bốn
người lập tức biến phòng tôi thành một nơi đầy ắp tiếng cười nói từ khi nào.

Hết chơi điện tử, họ lại cùng nhau kể chuyện đàm tiếu, giống như đám con gái học cấp
ba. Tôi thì không rõ thế nào, nhưng chắc chắn họ ngủ muộn hơn tôi.

Còn tôi? Tôi cũng chấp nhận chơi với họ vài trận, có thắng có thua, vui là chính. Còn về
``tiết mục'' đàm tiếu thì có lẽ tôi đã chịu hết nổi. Tôi đắp chăn, xoay mặt vào tường, quyết
định đi ngủ cho ``rảnh''.

\begin{center}
    \uline{Hiện tại.}
\end{center}

\dots Đúng, đêm hôm qua nhà tôi thực sự bát nháo như thế đấy.

Nhìn bên trên tấm nệm mà vốn dĩ tôi sẽ ngủ, tôi thấy Fubuki-chan và Mio-chan nằm sát
bên trái, còn Okayu-chan và Korone-chan thì chiếm lĩnh bên phải, ở tận bên trong của
tấm nệm.

Cáo trắng nằm nghiêng, tay ôm chặt một chiếc gối ôm như báu vật, khuôn mặt nở nụ cười
mơ hồ. Sói đen thì ngủ thẳng đơ, nhưng chiếc chăn tụt xuống tận eo, để lộ dáng ngủ kỳ
quặc mà\dots\ thậm chí tôi còn không dám nhìn lâu vì em ấy không mặc quần lót.

Mèo tím thì cuộn mình lại như một chú mèo béo ngủ đông, còn Chó nâu thì gác tay gác
chân thoải mái, đầu gối lên bụng bạn thân của em ấy, miệng còn khe khẽ ngáy. Con Mi
trắng của tôi, nằm gọn dưới chân tôi, nghe thấy tiếng động liền mở mắt, nhưng chỉ liếc
nhìn tôi vài giây rồi lại rúc vào chỗ ấm áp của nó.

Tôi quay đầu, thấy bố tôi đang ngồi dựa tường ở góc phòng, tay cầm cốc cà phê, mắt chăm
chú vào màn hình điện thoại. ``{\color{red} Bố lên từ khi nào vậy?}'' tôi khẽ hỏi.

Bố tôi ngẩng lên, trả lời bằng giọng đều đều: ``{\color{red} Vừa mới lên ban nãy. Cả đêm qua bố có
    ngủ được đâu.}''

Tôi chỉ im lặng, để ông ấy tiếp tục: ``{\color{red} Có chuyện gì mà đêm qua ở trên nhà nhốn nháo quá
    thế? Bà nội đêm qua cứ nhặng xị cả lên kia kìa.}''

Tôi nhún vai, ngáp dài: ``{\color{red} Con cũng chịu. Đêm qua lúc con ngủ thì thấy họ còn đang buôn
    dưa lê với nhau. Chắc là thức đến tận 2, 3 giờ sáng mới ngủ đây mà.}''

Bố tôi nhấp ngụm cà phê, nhếch môi cười: {\color{red} ``Đám bạn của con cũng loạn phết nhỉ.}'' Tôi
thở dài đáp: ``{\color{red} Cũng không hẳn là bạn đâu, chỉ là nhân viên với quản lý thôi ạ. Nhưng mà
    ở vế `loạn' thì con đồng ý.}''

Nói rồi, tôi chậm rãi rời khỏi đệm. Không khí lạnh buốt từ ngoài trời phả vào khiến tôi
rùng mình. Mùa Tết năm nay vẫn rét như mọi năm, đúng chuẩn kiểu Tết miền Bắc.

Tôi rửa mặt, đánh răng, rồi rót cốc cà phê còn dư từ phin bố pha ban nãy. Sau đó, lấy thêm
sữa đặc và vài viên đá để làm một ly bạc xỉu. Vừa cầm ly cà phê nóng vừa rùng mình vì
cái lạnh, tôi lại ngồi xuống cạnh bố, tận hưởng chút thời gian yên tĩnh hiếm hoi vào buổi
sáng.

\ct

Một lát sau, em gái tôi, Ryuuhou, cũng lục đục thức dậy. Tôi nhìn em kéo lê dép, gương
mặt mơ màng nhưng rồi khẽ bật cười khi thấy bốn người vẫn đang say ngủ trên đệm.

``{\color{red} Họ vẫn đang ngủ ạ?}'' em gái tôi hỏi, khẽ mỉm cười khi trông thấy khung cảnh bình yên
hiếm có từ bốn cô gái vốn dĩ siêu quậy này.

Tôi nhún vai: ``{\color{red} Ừ. Thử đoán xem họ ngủ lúc mấy giờ?}'' Cô em gái chống cằm suy nghĩ
vài giây, rồi đáp: ``{\color{red} Ờm\dots\ Ba giờ?}''

Tôi khẽ lắc đầu, rồi cất lên giọng xỉa xói em ấy: ``{\color{red}Đúng cái giờ mà em hay ngủ nếu hôm
    sau là ngày nghỉ đấy.}'' Ryuuhou khẽ bật cười.

Tôi lại hỏi: ``{\color{red} Mẹ vẫn chưa dậy à?}'' Em gái tôi ngừng cười, rồi đáp ngắn gọn một tiếng:
``Ừ,'' rồi ngồi xuống bên cạnh tôi, tiện tay lấy luôn cốc bạc xỉu mà bố vừa pha. Ba bố con
chúng tôi tiếp tục tận hưởng bầu không khí thư giãn này.

\ct

Tám giờ ba mươi sáng. Mẹ tôi cuối cùng cũng xuất hiện, dáng đi uể oải, vừa đi vừa ngáp
dài. Những hôm không đi làm, mẹ tôi thường ngủ đến trưa. Nhưng nhìn cách bà bước
lên cầu thang, tôi đoán chắc có chuyện gì đó bất thường khiến bà phải dậy sớm.

``{\color{red} Ở dưới kia ồn ào quá\dots}'' mẹ tôi lên tiếng, giọng đầy vẻ trách móc. ``Các bạn của con đến
sớm thật đấy.''

Tôi chỉ có thể nhíu mày và cười trừ: ``{\color{red}Giờ này chắc là giờ mọi người đi làm rồi mà.}''
Nhìn bốn người vẫn nằm lăn lóc trên đệm, mẹ tôi thở dài, lắc đầu ngán ngẩm: ``{\color{red} Thế này
    thì không biết bao giờ mới chịu dậy\dots}''

Tôi nói, chỉ có thể cười bất lực: ``{\color{red} Con nghĩ là mấy đứa đó biết rõ hôm nay là ngày nghỉ
    nên đến giờ vẫn còn nướng đấy ạ.}''

Bất chợt tôi nảy ra một ý tưởng: ``{\color{red} Con có trò này để trêu họ này,}'' rồi thì thầm với ba người
trong gia đình. Mỗi người đều có một phản ứng khác biệt:

Mẹ tôi tỏ vẻ hơi ngần ngại: ``{\color{red} Con chắc trò này không làm phiền các bạn của con chứ?}''

Tôi mỉm cười đáp: ``{\color{red} Không mẹ. Hoàn toàn vô hại ạ.}''

Bố tôi thì chỉ nhấp thêm ngụm cà phê: ``{\color{red} Đừng có làm quá là được,}'' rồi tiếp tục cắm đầu
vào đọc sách.

Riêng cô em gái tôi thì háo hức hơn hẳn, đến mức khẽ reo lên: ``{\color{red} Em theo phe anh luôn,
    Onii-san! Lâu lâu mới có dịp chơi lớn thế này!}''

Sau khi tôi đảm bảo rằng đây chỉ là một trò đùa nhỏ, không làm tổn thương ai, cả nhà
cũng miễn cưỡng đồng ý. Tôi nhanh chóng bày tỏ kế hoạch, khiến Ryuuhou cười khúc
khích đầy thích thú, còn mẹ tôi thì thở dài, nhắc khéo: ``{\color{red} Xong rồi phải xin lỗi đàng hoàng
    đấy nhé.}''

\ct

Tôi chạy xuống tầng 3, nơi tất cả các thành viên (trừ A-san) đã tụ tập.

Miko-chan đang lục tục chỉnh lại chiếc dây buộc bụng, trong khi Matsuri-chan và Akutan
cãi nhau về chuyện tối qua ai là người ăn hết bánh kẹo trong bếp. Subaru-chan ngồi phịch
xuống ghế, tay xoa xoa vai như thể vừa luyện tập gì đó căng thẳng lắm. Sora-chan thì dịu
dàng rót trà cho Mel-chan và Ayame-chan, còn Peko-chan thì đang nghịch nghịch tai thỏ
của chính mình.

Haato-chan chợt lên tiếng: ``Không biết cái đám \textit{GAMERS} đi đâu nhỉ\dots'' Nàng Thỏ gật
đầu phụ họa: ``Tối qua bảo là đi có việc mà nhỉ? Sao bây giờ vẫn chưa thấy về\dots''

Noel-chan cau mày: ``Cơ mà lúc đi, thấy mặt của Fubuki-senpai, Okayu-senpai với Koronesenpai
trông nghi lắm\dots'' Mel-chan tiếp lời, giọng đầy băn khoăn: ``Nhưng mà Mio-chan
cũng đi nữa\dots\ Không lẽ\dots''

Roboco-chan thì phán đoán với vẻ hồn nhiên như mọi khi: ``Chị nghĩ là họ sang nhà
Kuon-senpai ngủ sao\~{}''

Ngay khi em ấy vừa nói xong, tôi bật cười: ``Khá khen cho một đứa PON như em lại đoán
đúng đấy.'' Như một phản ứng theo lẽ tự nhiên, em ấy đã giãu nảy: ``Đã nói rồi, em không
có PON mà!''

Tuy nhiên, lời khẳng định của tôi khiến cho cả căn phòng bỗng rơi vào im lặng trong vài
giây. Vẻ mặt ai nấy đều tỏ vẻ ngạc nhiên đến tột độ. Sora-chan thậm chí đang rót dở trà
còn tràn hết ra ngoài.

Và rồi, ba, hai, một\dots

``\textbf{HẢ!!??}''

Ngay lập tức, tất cả mọi người xúm lại quanh tôi, hối hả hỏi cho ra lẽ. Choco-sensei
nghiêng người, giọng đầy hoài nghi: ``Có đúng không vậy, Kuon-sama?''

Tôi gật đầu: ``Đúng mà.'' Ojou-san cau mày: ``Anh mời họ lên sao?'' Tôi nhún vai: ``Họ
tự lên thôi. Cả đêm hôm qua anh có ngủ được đâu.''

Đó rõ ràng là lời nói dối. Nhưng nói thật làm gì, khi thêm chút kịch tính thế này thú vị
hơn nhiều.

Thấy rằng lời nói của tôi có phần hơi mâu thuẫn, Subaru-chan thắc mắc: ``Em thấy có gì
đó sai sai\dots\ Nếu mà anh không ngủ được thì tại sao anh lại tỉnh như sáo thế?-''

Tôi tỉnh bơ: ``Anh vừa uống cà phê mà. Uống phát tỉnh luôn.'' Đó cũng rõ ràng là lời nói
dối. Phần lớn là dối. Đôi khi tôi uống hết cốc cà phê, rồi cuối cùng vẫn lăn đùng ra ngủ
như bình thường.

Sora-chan nghiêng đầu: ``Nhưng mà\dots\ nhỡ đâu bạn của anh đến chơi thì sao?'' Rushiachan
run rẩy, như thể em ấy đang nghĩ rằng tôi đang nói đùa: ``Đ-Đúng! Có khi anh ấy đi
dọa chúng ta thì sao nhỉ?''

Thấy rằng mọi người vẫn chưa tin lời mình, tôi rút điện thoại ra, giơ lên trước mặt mọi
người: ``Đây này, bằng chứng rõ ràng.''

Tôi mở album ảnh đêm qua ra, lần lượt cho mọi người xem. Bốn bức ảnh đầu tiên là ảnh
nhóm \textit{GAMERS} ``tạo dáng'' đủ kiểu, còn bức cuối là cảnh cả nhóm ngồi tám chuyện rôm
rả cùng với tôi ở góc dưới bên phải. Tôi chụp bức ảnh đó một cách rất bâng quơ ngay
trước khi ngủ, không ngờ nó sẽ lại hữu ích.

Ru-chan thấy vậy, nhanh như chớp giật lấy điện thoại, lướt qua từng bức ảnh với tốc độ
ánh sáng. Các cô gái khác ngay lập tức ùa vào xem, xô đẩy nhau chen chúc. Ngay khi em
ấy dừng lại ở bức ảnh cuối, sắc mặt nàng Cô đồng tối sầm. Tay em ấy bắt đầu run lên, rồi
bất chợt siết chặt cái điện thoại như muốn bóp nát nó.

``N-Này! Bình tĩnh nào!'' tôi vội hét lên, nhưng chẳng giúp ích gì. Ru-chan gằn giọng,
đôi mắt tóe lửa: ``\textbf{Sao bọn bây dám!?}'' rồi định rút con dao phép ra, lên phòng tôi tính
``xử đẹp'' bốn người bọn họ.

``Chết rồi!'' Tôi, Mel-chan và Sora-chan phải lao vào ngăn cản. Mất cả một hồi vật lộn
mới có thể giữ được Rushia, để điện thoại không bị vứt đi hoặc ai đó không bị đâm.

Nhìn cái cảnh hỗn loạn trước mắt, tôi chỉ có thể thầm thở dài: ``\textit{Mình đi hơi xa thật rồi\dots}''
Nhưng dù gì, chúng tôi vẫn sẽ phải tiếp tục với kế hoạch này.

\ct

Dù rằng Cô đồng đã bình tĩnh lại, không khí trong phòng vẫn còn căng thẳng. Một vài
người, đặc biệt là Rushia-chan, tỏ rõ sự bực bội, trong khi những người khác lại xen lẫn
chút ghen tị.

Shion-chan khoanh tay, bĩu môi: ``Thật không thể chấp nhận được!'' Aqua-chan cũng phụ
họa: ``Đúng đó, bất công thật sự!''

Miko-chi phồng má, giọng pha chút bất mãn: ``Mà tính ra, mấy người đó đúng là thâm
đấy ne\dots'' Peko-chan giật tai thỏ, lẩm bẩm: ``Định qua mặt tụi này à, peko\dots''

Bản thân Ru-chan nghiếng răng ken két, ánh mắt đầy hận thù: ``Vị trí đó phải là của em
mới đúng!''

Nhìn khuôn mặt đỏ bừng của mọi người, tôi cố nhịn cười, nhưng đồng thời cũng hiểu rằng
nếu không can thiệp, cơn bão ``trả thù'' sẽ bùng lên bất cứ lúc nào.

Tôi liền giơ tay trấn an: ``Rồi, rồi, cứ bình tĩnh nào. Anh biết là mọi người đang tức giận.
Đột nhiên bị bán đứng chỉ vì muốn ngủ qua đêm với một người con trai, đúng không?''

Tất cả im lặng, khuôn mặt của họ đỏ ửng. Một vài người lúng túng liếc nhau, nhưng cuối
cùng cũng gật đầu đồng tình.

Một số người lên tiếng bức xúc, tỏ vẻ sự bực bội vì bị đồng đội trong nhóm ``phản bội.''
Tỏi-chan hậm hực: ``Chứ còn sao nữa!''

Hành tím-chan thì bĩu môi: ``Thật là không thể chấp nhận được!'' Ru-chan thậm chí còn
không thể kìm nén được cơn giận, bừng bừng giận dữ: ``Làm sao mà em chịu được chuyện
này chứ!''

Thấy em ấy chuẩn bị ``nổi trận lôi đình'' lần nữa, tôi vội nói: ``Rồi, rồi. Anh biết. Mấy
đứa muốn trả thù đúng không?''

Tất cả mọi người cùng đồng thanh: ``\textbf{Quá đúng!}'' Tôi nở nụ cười đầy bí hiểm: ``Vậy thì
nghe anh nói đây. Kế hoạch là thế này\dots''

Sau khi tôi giải thích kế hoạch, cả nhóm đều đồng ý một cách hào hứng. Từ ánh mắt rạng
rỡ đến những cái gật đầu chắc nịch, ai cũng tỏ ra rất thích thú với màn ``trả thù'' sắp tới.

Nàng Y tá phấn khích: ``Trò này vui đây!'' Nàng Lễ hội bật cười lớn: ``Đúng gu của em
luôn!''

Nàng Phù thủy vỗ tay, hối thúc mọi người: ``Được rồi, làm thôi!''

\ct

Chúng tôi - bao gồm tôi, Choco-sensei, Mel-chan, Haato-chan, Shion-chan, Sora-chan,
Matsuri-chan và Ryuuhou - cùng rón rén lên tầng năm, nơi mà các cô nàng \textit{GAMERS} vẫn
đang say giấc. Tôi nhẹ nhàng mở cửa phòng, bước vào.

Bên trong, bốn cô gái nằm ngủ không chút cảnh giác. Fubu-chan cuộn mình trong chăn,
còn Mio-chan nằm nghiêng với tay vắt qua đầu, chiếc quần lót vẫn đang vứt tứ tung bên
cạnh con Mi đã tỉnh giấc từ khi nào. Okayu-chan nằm úp, mặt đè xuống gối, trong khi
Koro-chan thì\dots\ xoay người lung tung, chân đá lên không trung.

Tôi khẽ lay từng người: ``Mấy đứa, sáng rồi mà vẫn chưa dậy à?''

Phản ứng đúng như tôi dự đoán. Cáo trắng chỉ khẽ nhíu mày rồi đổi tư thế nằm. Sói đen
lẩm bẩm điều gì đó rồi kéo chăn trùm kín hơn. Mèo tím thì không có chút động tĩnh, còn
Chó nâu\dots\ bất thình lình đá tôi một cú khiến tôi loạng choạng suýt ngã.

Thấy cả bốn cô gái vẫn không có ý định ngủ dậy, tôi cười khì: ``Đúng kiểu.'' Rồi tôi quay
sang cả nhóm thì thầm: ``\hl{Bắt đầu triển khai!}''

\ct

Chúng tôi nhẹ nhàng khiêng từng người xuống tầng 3, thật may mắn là không ai tỉnh giấc
giữa chừng. Điều đáng kinh ngạc là dù sở hữu tai chó, tai mèo thính như vậy, họ vẫn
không để ý chuyện gì đang xảy ra với họ.

Trong đó, tôi và nàng Cổ động khiêng Fubu-chan, nàng Dơi và nàng Phù thủy thì vác
nàng Sói đen xuống (với quần lót được Ryuuhou mang xuống nhà để đem đi giặt), nàng Song tính và nàng Thần tượng thì cõng cặp đôi OkaKoro xuống, còn Sensei thì làm ``xi
nhan'', hướng dẫn chúng tôi đi xuống.

Khi đã vác bốn người xuống tầng ba, chúng tôi đặt họ lên bốn tấm thanh gỗ lớn lấy ra từ
kho. Với sự phối hợp ăn ý, cả nhóm buộc tay chân từng người vào thanh gỗ, cố gắng làm
thật nhẹ nhàng để không làm họ tỉnh.

Một lát sau, khung cảnh hoàn tất. Bốn cô nàng \textit{GAMERS} bị buộc chặt trên những thanh
gỗ, trông chẳng khác nào các ``tội phạm'' thời Trung Cổ đang chuẩn bị cho buổi xử án.
Chỉ khác ở chỗ, chẳng có ai chết, cũng chẳng có ai bị đau đớn.

Khi xong xuôi, tôi bước lùi lại vài bước, giơ tay ra hiệu cho cả nhóm. Rồi, mọi người lùi
hết về vị trí đi. Anh sẽ chạy lên nhà. Bắt đầu phần chính!''

Tôi nhanh chóng chạy lên tầng trên, chuẩn bị quan sát toàn bộ màn kịch từ xa. Pekorachan,
với nụ cười đầy tinh quái, bắt đầu đếm: ``3\dots\ 2\dots\ 1\dots\ Peko!''

\pov{-}

Nàng Thỏ vừa dứt lời, hai cái xô nước lớn mà nàng Roboco và Noel chuẩn bị sẵn từ trước
liền được đổ thẳng xuống, dội toàn bộ nước lạnh lên bốn cô gái \textit{GAMERS}.

``Gyaaaah!!'' Một tiếng hét đồng thanh vang lên khi cả nhóm bị đánh thức đột ngột.
Fubuki và Mio, với bản năng nhanh nhạy, bật dậy đầu tiên, trong khi Okayu và Korone
vẫn còn mơ màng, gương mặt lộ rõ sự mệt mỏi.

``C-Có chuyện gì vậy!?'' nàng Cáo trắng vừa hất nước khỏi mặt, vừa lắp bắp hỏi, ngơ
ngác nhìn mọi người. Nàng Sói đen thì nhìn quanh, nhưng chỉ kịp nhận ra tay chân mình
đang bị trói chặt từ phía sau. Cô hoảng hốt thốt lên: ``S-Sao tay chân của tôi lại bị trói thế
này!?''

Đáp lại sự hoang mang của bốn người, giọng nói lạnh lùng của Shion vang lên từ phía xa:
``Bốn người tỉnh chưa?''

Bốn cô gái quay đầu về phía âm thanh, thấy Shion đang đứng chống hông với vẻ mặt
không mấy hài lòng. Xung quanh, tất cả các thành viên khác của \hll\ (trừ A-chan, vì
cô chưa có mặt) đang đứng đó, mỗi người một vẻ: người buồn rầu, người bất mãn, đặc
biệt là ánh mắt bừng bừng sát khí của Rushia.

Trưởng nhóm \textit{GAMERS} nuốt khan, cố gắng giữ bình tĩnh: ``S-Sao mọi người lại ở đây?''
Pekora, với nụ cười nửa miệng, bước lên phía trước, đặt câu hỏi đầy ẩn ý: ``Khai thật đi,
peko. Tối hôm qua bốn người đã đi đâu?''

Korone, dù vẫn còn ngái ngủ, cố gắng trả lời: ``Bọn chị đi chơi thôi mà\dots'' Haato nhíu
mày, xen ngang: ``Ở nhà Kuon-senpai à?''

Nàng Cáo trắng vội ngẩng mặt lên, vội vàng phân bua: ``Cái đó chỉ là bọn tôi ghé qua\dots\ một chút thôi\dots''

Mel nheo mắt: ``Thế sao tôi lại nghe nói là bốn người đã ngủ cùng với anh ấy nhỉ?'' Câu
nói của nàng Dơi khiến cả bốn cô gái bị trói trên thanh gỗ sững sờ. Cáo trắng lắp bắp: ``C-Cái đó thì\dots''

Không để họ có cơ hội bào chữa, Shion hất cằm, giọng đầy tự tin: ``Đừng chối. Em đã
theo dõi cả bốn người tối hôm qua rồi. Nhìn đi.'' Cô vung tay, sử dụng phép thuật của mình để chiếu
lên không trung năm bức ảnh từ điện thoại của Kuon. Hình ảnh rõ mồn một cảnh Fubuki, Mio, Okayu và
Korone nằm ngủ say sưa trong phòng của cậu con trai.

Toàn bộ nhóm \textit{GAMERS} đứng hình, trong khi những người còn lại ồ lên ngạc nhiên. Rushia còn hét lớn:
``Đó! Chính họ đấy!''

Nàng Sói đen vội vàng thanh minh trước mặt mọi người: ``Không phải đâu! Tại Fubuki cứ nằng nặc rủ
tôi đi chứ bộ!'', đồng thời quay sáng người bạn vẫn đang ngỡ ngàng kia: ``Đó, thấy chưa! Tôi đã bảo
đây là ý tưởng tồi mà!''

Trong khi đó, nàng Mèo tím ngáp dài, đôi mắt vẫn còn lờ đờ, công khai một cách thẳng thắn: ``Em thì\dots\ tại Koro-san mời đi thôi mà\dots''

Nàng Khuyển trố mắt nhìn bạn của cô, mồ hôi rịn trên trán: ``Ogayu, bà định phản bội tôi sao!? Cái
này là do Fubuki bày trò mà!''

Nàng Cáo trắng, chưa kịp phản bác, đã bị nàng Thỏ cắt ngang bằng giọng nghiêm nghị: ``Đủ rồi, peko.
Chúng tôi không quan tâm ai bày trò trước.''

Choco lắc đầu, giọng đầy thấy vọng: ``Mọi người cứ đi trước như thế này\dots\ Cô thực sự thất vọng với
các em quá\dots'' Haato thì bĩu môi bất bình: ``Mấy người chơi thế thì ai chơi lại được!?''

Miko thì nắm chặt tay, nhìn họ với ánh mắt kiên quyết: ``Nếu cứ để thế này thì chúng ta
sẽ chằng còn cơ hội gì nữa, ne! Vì vậy\dots''

Cuối cùng, Rushia, người từ nãy giờ chỉ đứng yên nhìn, bước lên phía trước. Nàng tóc
xanh nước biển khoanh tay, ra phán quyết như một vị Bao Công: ``Cả bốn người phải chịu
phạt!''

Lời tuyên bố của Rushia khiến bốn cô nàng \textit{GAMERS} hốt hoảng, đồng thời làm những
người xung quanh bật cười thích thú, thậm chí còn hò reo ủng hộ. Ai nấy đều hào hứng
chờ xem ``hình phạt'' sắp tới sẽ như thế nào.

Nàng Thỏ tiếp lời cô, giọng toát lên vẻ uy quyền: ``Mio-senpai và Okayu-chan, vì hai
người chỉ là những kẻ bị kéo theo, nên sẽ nhận hình phạt nhẹ hơn và bị hành quyết trước.''

Mio nhíu mày, không giấu được vẻ hoang mang: ``Lại còn có hình phạt nhẹ với nặng nữa
à? Rồi trước sau là sao? Hành quyết kiểu gì kỳ lạ vậy?'' còn Okayu thì ngáp nhẹ, hỏi với
giọng tỉnh rụi: ``Vậy hình phạt đó là gì thế?''

Đáp lại câu hỏi của hai người, tất cả những thành viên còn lại đồng thanh hét lên: ``\textbf{Xịt kem!}''

Chưa kịp hiểu chuyện gì, hai nàng đã bị cả đám cầm bình xịt kem trắng mà Shion chuẩn
bị trước phun thẳng vào người.

Nàng Sói đen hét toáng lên: ``Á!! Sao nhiều quá vậy!?'' Nàng Mèo thì cười khúc khích,
bất chấp tình cảnh của mình: ``Ui, cái này làm em thấy hơi buồn (nhột) nha\~{}''

Korone, người đang bị cắm cọc ngay bên cạnh, trố mắt nhìn cô: ``Ơ\dots\ Ogayu cười vì bị
kem chọc lét à?'' Fubuki thì cười với vẻ thích thú: ``Trông có vẻ vui à nha\~{}''

Chẳng mấy chốc, toàn thân hai người bị phủ đầy kem trắng xóa. Trong lúc cao hứng,
Noel vô tình bóp quá tay, khiến cả bình xịt của cô bay thẳng về phía trước, phun bọt kem
khắp phòng.

Nàng Hiệp sĩ hốt hoảng: ``Thôi chết\dots\ Em lỡ ta rồi!'' Roboco đứng kế bên cũng cuống
quýt không kém: ``Chị cũng thế\dots\ Ôi không!''

Cả nàng Người máy cũng bóp quá tay, và bình xịt kem của cô cũng bắt đầu bay tứ tung.
Cả hai bình cứ bay quanh khắp phòng, và rồi *KENG!* Chúng vô tình cụng vào nhau
và bay thẳng đến vị trí của Fubuki và Korone. Cả hai chỉ kịp thốt lên một tiếng ``Ah\dots''
trước khi bị bình xịt tông thẳng vào bụng dưới. Tiếng ``\textbf{Hự\dots!}'' vang lên rõ to, cả hai run
lẩy bẩy trên thanh gỗ. Riêng nàng Khuyển thì không chịu nổi mà ``dấm đài'' ngay tại chỗ.

Miko, đứng ở dưới, bật cười khúc khích: ``Nhìn giống vụ Aki-chan với Mio-chan bị bím
tóc đập trúng hồi trước\footnote{Xem {\flaregothic ÉPISODE 20}, quyển số Hai.} ghê ne\dots\ Đau lắm đấy.''

Cả hai người khi bị cô nàng réo tên đều đồng loạt nói, mặt của họ có chút ửng đỏ: ``\textbf{Tất
    nhiên là đau rồi!}''

Pekora cũng không giấu được sự ngạc nhiên: ``À-rế? Cái này không có trong dự tính đâu
nhé, peko.'' Shion chép miệng: ``Cũng tại hai người hậu đậu đó thôi.'' Mel nhìn cả hai
nạn nhân bị tông kia, thở dài: ``Nhưng trông họ đau thật đấy\dots\ Nhất là Korone-chan\dots''

Haato lên tiếng, phá vỡ sự im lặng: ``Thôi thì hình phạt cho Mio-senpai và Okayu-chan
xong rồi, còn hai người kia thì sao?''

Nàng Thỏ xanh sau một hồi suy nghĩ, bèn đáp: ``Ta sẽ cù lét họ, peko!'' Nghe thấy vậy,
không ai bảo ai, họ lập tức lôi ra những chiếc chổi lông gà mà họ ``xin'' được từ kho nhà
cậu Tổng.

Matsuri nhíu mày với công cụ ``tra tấn'' họ đã chuẩn bị trước: ``Ờm\dots\ Ta sẽ cù lét bằng
mấy cái này á?'' Sora cũng gật gù, góp ý với mọi người: ``Dùng tay không chẳng phải
tiện hơn sao?''

Mọi người nhìn nhau, đồng tình với lời nói của nàng Thần tượng. Họ bỏ hết chổi lông gà
xuống, quay sang bốn cô gái rồi giơ hai bàn tay lên, lắc lắc ngón tay.

Fubuki hoảng hốt: ``A-A-nô\dots\ Mọi người định làm gì vậy?'' Rushia cười nham hiểm (dù
không còn nguy hiểm như hồi ban nãy): ``Coi như đây là cách bọn này trả thù đi-nano
desu.''

Mio kêu lên, cố cứu vãn tình thế cho người bạ của cô, dù vẫn còn khá bất bình: ``Thôi
nào, tha cho cậu ấy đi mọi người.''

Riêng Korone, mặc dù mặt tái xanh, nhưng vẫn cố tỏ ra nguy hiểm: ``Ngon thì nhào vô!''
Điều này làm nàng Sói đen kinh hãi quay sang: ``K-Korone!?''

Nàng Khuyển cười đắc ý: ``Dù gì em cũng là `God Dog' mà, em chịu được hết!'' Thấy
thế, Ookayu liếc xéo bạn mình, khịa: ``Bà giống Hot Dog vừa bị tè dầm thì đúng hơn á\~{}''

Câu nói như xát vào nỗi xấu hổ của cô khiến cho cô giãy giụa trên thấm thanh gỗ, mặt đỏ
bừng: ``Là `God Dog'! Với lại tôi chỉ bị bình xịt tông trúng thôi!''

Đó cũng là câu cuối cùng hai nàng kịp thốt lên trước khi toàn bộ gần hai mươi người còn
lại lao đến, bắt đầu cúi người cù lét khắp toàn thân của cả bốn người.

Nàng Cáo trắng gào lên: ``\textbf{D-Dừng lại đi- Ahahaha!}'' Nàng Miêu tím thì vừa cười vừa
lẩm bẩm: ``Nữa đi, nữa đi\~{}''

Nàng Chó nâu thì cố nhịn, nhưng rồi cuối cùng cũng bật ra: ``Mình chịu được, mình chịu
được, mình- Ahaha!'' còn nàng Sói đen chỉ biết kêu trời: ``Tại sao tôi cũng bị dính chưởng
chứ- Ahahaha!''

Tiếng cười vang lên như sóng cuộn, kéo dài suốt hơn hai mươi phút. Đến cuối cùng, bốn
người trên thanh gỗ, kiệt sức, đồng loạt cầu xin: ``Bọn em/tôi hứa sẽ không lừa dối mọi
người nữa mà!''

Càng tệ hơn, toàn bộ sự việc, từ lúc đưa họ xuống đến khi họ cầu xin, đã được Rushia và
Shion quay lại bằng máy quay mà họ mượn từ A-chan.

\ct

\pov{Hikawa Kuon}

Đúng lúc này, tôi nhìn thấy A-san xuất hiện, tiến về căn phòng ở tầng ba, nơi mà chúng
tôi sẽ tụ tập lại vào những ngày sắp tới. Chị ấy đi với dáng vẻ điềm tĩnh, miệng khóe một
nụ cười khi chị nghĩ rằng hôm nay sẽ là một ngày thích hợp để tiếp tục dự án quay phim.

Thế nhưng, đập vào mắt chị là một tình cảnh hỗn loạn trước mặt. ``Mọi người lên đây
sớm thế\dots\ Có chuyện gì vậy?'' chị ấy hỏi, giọng pha chút khó hiểu.

Tôi vừa đặt chân xuống tầng cũng phải ngẩn người: ``Mấy đứa làm đến đâu rồi- Âu sịt\dots''
Mắt tôi chạm ngay đống bọt kem trắng xóa lan tràn khắp phòng. Các thành viên nhóm
GAMERS, vẫn bị treo trên thanh gỗ, thở hồng hộc, thân mình lẩy bẩy sau những phút tra
tấn cười.

``\textbf{Thả bọn tôi ra!}'' Mio-chan hét to, giọng đầy uất ức, nhưng chẳng ai trong đám còn lại
chịu nhúc nhích. Thậm chí tôi còn thấy họ còn chọc tức bốn cô gái ấy nhiều hơn.

A-san quay sang tôi, vẻ dò xét: ``Sao nhà cậu lại tứ tung bọt kem thế này?'' Tôi giả vờ
ngơ ngác, khoanh tay lắc đầu: ``Em đâu có biết gì đâu, em vừa mới lên mà,'' câu trả lời
kèm theo nụ cười mỉm ``vô tội.''

``Sao cũng được,'' chị thở dài. ``Mau phụ tôi đưa họ xuống đi.''

Chúng tôi tiến vào phòng rồi thả từng người trên thanh gỗ xuống. Mio-chan, Okayu-chan, Korone-chan và Fubuki-chan vừa được giải thoát liền đổ sụp xuống sàn, thở dốc như thể
vừa chạy ma-ra-tông.

Cả bốn người đều nhìn tôi với ánh mắt dò xét, như thể họ đang tỏ vẻ oan ức. Tôi chỉ có
thể cười trừ nhìn những cặp mắt đang rơi lệ đó, thầm nói: ``\hl{Xin lỗi mấy đứa nha\dots}''

Đúng lúc đó, từ phía xa, hai cô nàng tinh quái là Peko-chan và Ru-chan bước ra, miệng
reo vang: ``Nhiệm vụ trả thù hoàn tất!''

A-san đẩy gọng kính, giọng nghiêm nghị: ``Tôi muốn mọi người giải thích, chuyện gì vừa
xảy ra ở đây thế này?''

Và thế là, từng chút một, toàn bộ sự việc được kể lại: từ ý tưởng trả thù cho đến màn
xịt kem ``đầy sáng tạo,'' thậm chí là pha ``cụng bình'' khó đỡ ấy. Điều đáng ngạc nhiên là
không ai nhắc đến vai trò của tôi trong vụ này, có lẽ vì tôi đã giúp họ khá nhiều việc trước
đó.

Chị ấy lắng nghe chăm chú, cuối cùng quay sang bốn nạn nhân bất đắc dĩ. ``Vậy\dots\ Shirakami-san, Ookami-san, Nekomata-san và Inugami-san tự ý rời khỏi khu tập thể để
sang nhà Hikawa-san ngủ, mà không có sự đồng ý của tôi, đúng không?''

Sói đen-chan cố gắng chống chế: ``Tại Fubuki muốn rủ em đi cùng chứ đâu em có muốn!''
Cáo trắng-chan thì vội phản bác ngay: ``Tôi đâu có ngờ mọi chuyện sẽ thành ra như thế
này đâu!''

Thấy hai người chuẩn bị cãi cọ nhau, A-san giơ tay ra hiệu cho họ im lặng: ``Đủ rồi!'' và
quay sang hỏi tôi: ``Hikawa-san, có đúng là họ sang nhà cậu ngủ không?''

Tôi gật đầu, giọng ôn tồn: ``Có chị ạ. Họ cứ thế sang thôi, không báo trước với em một
tiếng.'' Đó thực sự là một lời nói thật, không ai bảo tôi rằng bốn người đó sẽ sang nhà tôi
ngủ cả!

Chị ấy thở dài, như đang cố kiềm chế sự khó chịu: ``Các cô có biết là không được tự ý
ngủ chung với quản lý, đúng chứ?''

Mio-chan ấp úng: ``Cái đó thì em biết, nhưng\dots''

``Tôi không cấm mọi người ngủ với Hikawa-san,'' chị cắt ngang, tiếp tục nói với giọng
đanh thép, ``nhưng ít nhất phải báo trước với tôi và cậu ấy một tiếng chứ!''

Okayu-chan, vốn ít khi lúng túng, cũng buộc phải lên tiếng: ``Bọn em\dots\ thấy chị làm việc
hăng quá nên không dám làm phiền ạ.''

Koro-chan gật gù thêm vào: ``Đúng á, A-chan làm việc chăm chỉ quá mà.''

Mặc cho cặp đôi đưa ra lý do khá thuyết phục (ít nhất là với tôi), chị vẫn giữ thái độ cứng
rắn: ``Biết là thế, nhưng mà, hành vi của các cô có thể gây ảnh hưởng tới gia đình cậu ấy
nữa. Các cô có nghĩ đến chuyện đó không?''

Fubu-chan vội vã đáp lại: ``Nhưng anh ấy nói là không phiền mà!'' Lập tức, em ấy bị
A-san vặn lại, nhấn mạnh: ``Nhưng nếu gia đình cậu ấy thấy phiền thì sao?''

Koro-chan đáp ngay: ``Bố mẹ anh ấy thoải mái với chuyện đó mà chị\dots'' Đến đây, A-san
dừng lại, hít một hơi sâu. ``Có thể gia đình họ không ngại, nhưng không phải ai cũng thế
đâu! Thế nên\dots''

Và thế là, màn ``giảng đạo'' dài đến ba phút đồng hồ bắt đầu. Nếu so với việc tôi mắng
xa xả Robo-PON hơn 15 phút vì phá sập cả cái công ty hồi ba tháng trước, ít nhất họ vẫn
còn nhẹ chán.

Những thành viên còn lại, vốn vẫn đang nấp sau đống bọt kem, không khỏi cảm thấy hả
hê trước cảnh tượng này. Bản thân tôi thì chỉ biết cười trừ, cảm thấy có phần hơi tội cho
họ vì vướng phải mớ bòng bong này.

Tất nhiên, các cô gái ``tra tấn'' bốn cô gái \textit{GAMERS} đó cũng không thoát khỏi được tội,
khi A-san bất ngờ quay sang đống kem đó và chỉ thẳng tay vào đám đông: ``Cả các cô
nữa! Hạ nhục người ta rồi clay clip cũng chả hay ho gì đâu mà cười!''

Shion-chan từ đống bọt kem bước ra, giọng lí nhí: ``Bọn em xin lỗi chị, nhưng mà\dots'',
nhưng Ru-chan lại không tỏ vẻ chút ăn năn gì: ``Được trả thù thế này thoải mái lắm-nano
desu.''

\dots Thật sự thì Rushia-chan đang có một chút vấn đề về tâm lý đấy.

A-san nghe hai cô nàng tinh quái nói vậy, chỉ còn biết thở dài bất lực. Tôi tiến tới, nhẹ
nhàng đặt tay lên vai chị: ``Thôi chị ạ. Đằng nào hôm nay cũng cuối năm, có bao nhiêu
bực dọc cứ trút hết đi, không sang năm mới lại khổ.''

Mio-chan cũng gật đầu đồng tình: ``Dù khá uất ức với hành động của mọi người nhưng
em cũng phải công nhận anh ấy nói đúng.''

Mel-chan cười nhẹ: ``Nếu hôm nay là ngày cuối năm thì tại sao ta không ngồi lại và giải
quyết mọi bức bối trong lòng nhỉ?''

``Có thể chúng ta sẽ có cách giải quyết nào đó mà mọi người muốn tâm sự,'' Choco-san
tiếp lời, mỉm cười và hướng về phía tôi, ``Ví dụ kiểu Kuon-sama có thể chữa trị chứng
khó ngủ bằng một buổi ASMR chẳng hạn.''

Nhìn thấy em ấy nháy mắt, tôi đổ mồ hôi hột: ``Ờm\dots\ Cảm ơn, nhưng mà anh không cần
đâu.''

Câu nói của em ấy khiến cả nhóm cười phá lên, phá vỡ bầu không khí nghiêm túc ban
nãy. Tôi chỉ biết gượng gạo lắc đầu, cố gắng chuyển chủ đề trước khi mọi người làm quá
thêm nữa.

\ct

Vậy là, sau vài phút đùa cợt, tất cả cùng nhau xắn tay áo dọn dẹp tầng ba. Những lớp bọt
kem trắng xóa được lau sạch, các mảnh vụn và đồ đạc bị xô lệch được trả về vị trí cũ. Tôi
không ngờ rằng với một đám hỗn loạn như vậy mà họ vẫn phối hợp với nhau khá hiệu
quả.

Khi công việc hoàn tất, cả nhóm ngồi lại quanh một chiếc bàn lớn mà A-san đã chuẩn bị từ trước. Những câu chuyện được mở ra như dòng chảy bất tận: có chuyện tréo ngoe
khiến ai nấy đều ôm bụng cười, có những tâm sự thật lòng mà bấy lâu nay chẳng ai đủ
can đảm nói ra. Tôi có thể kể một vài điều ở đây.

Sora-chan kể về lần tập nhảy đến kiệt sức nhưng vẫn không làm được như ý muốn, chỉ
để rồi nhận ra rằng không ai yêu cầu em ấy phải hoàn hảo cả. Ayame-chan nhắc lại một
trò đùa ngớ ngẩn mà cô từng bày ra để trêu Mio-chan, kết quả là bị chính ``con Sói đen''
phục thù bằng cách vẽ thêm ``ria mèo'' lên mặt khi cô ngủ quên trong phòng tập.

Hay khi Mel-chan và Aki-chan chia sẻ về những ngày đầu bước chân vào \hll, khi
mọi thứ đều mới mẻ và khó khăn. Giọng nói của họ nhẹ nhàng nhưng chứa đầy sự trân
trọng dành cho khoảng thời gian đã qua.

Trong lúc chúng tôi đang bàn chuyện với nhau vừa nghiêm túc, vừa vui vẻ, tôi nhìn thấy
Sora-chan nhìn xung quanh căn phòng. Quả đúng là căn phòng ấy đã được dọn dẹp sạch
sẽ, nhưng có vẻ nó vẫn còn thiếu thốn gì đó\dots

Bất chợt, em ấy cất giọng hỏi: ``Kuon-senpai, Tết mà không có cành đào thì thiếu lắm
đúng không?''

``Ừ. Thảo nào anh cứ thấy thiếu thiếu,'' tôi đáp, khẽ gật gù đồng tình với em ấy. Cành đáo
ngày Tết là một trong những thứ không thể nào không có trong những dịp này được, khi
nó tượng trưng cho sự may mắn, niềm vui và hy vọng cho một năm phía trước.

Nhưng, cách mà chúng tôi mua cành đào lại\dots\ khá dị, khi ngay lúc chúng tôi vừa nhắc
đến cành đào xong, Shion-chan bằng một cách thần kỳ nào đó đã triệu ra một chậu đào
lớn, cành lá sum suê, còn treo thêm cả dây lục lạc màu đỏ rực rỡ.

Tôi ngỡ ngàng, hỏi: ``Shion-chan, em lấy cái này ở đâu vậy?''

À thì\dots\ một cửa hàng ở \ruby{Nisshin}{Nhật Tân}!'' Nàng Pháp sư nhí nhảnh đáp, nhưng tôi có cảm giác
câu trả lời này còn đáng ngờ hơn cả hành động của cô nàng.

Không dừng lại ở đó, Okayu và Korone còn lôi ra vài chậu quất vàng ươm, khiến căn
phòng tầng ba trông như một góc Tết thực thụ. Thực sự, không biết các cô gái đã ``nghiên
cứu'' và nghiền ngẫm khá đọng về văn hóa Tết ở Etsunan từ lúc họ tới đây.

Buổi sáng khép lại trong không khí vui vẻ và ấm áp. Khi kim đồng hồ chỉ đến trưa, mọi
người lục tục kéo nhau về nhà để dùng bữa và nghỉ ngơi.

Buổi chiều hôm đó diễn ra thật yên bình. Sau những gì vừa xảy ra, tôi chỉ mong khoảng
lặng này sẽ kéo dài thêm một chút nữa trước khi những ``rắc rối'' khác tìm đến.

\ct

Đã sáu giờ tối. Còn đúng sáu tiếng nữa là bước sang năm mới. Phố xá trước cửa nhà tôi
dần trở nên yên ắng, đến mức có cảm giác như mọi người đều đồng loạt rút về tổ ấm để
chuẩn bị đón giao thừa.

Thật ra, sự vắng vẻ ấy không hẳn là tự nhiên. Matsuri-chan và Noel-chan đã nghịch ngợm
mở hai cổng không gian từ trò chơi P*rtal (tôi vẫn không biết họ lấy chúng ở đâu) đặt tại hai đầu đường. Nhờ thế, cả khu phố trở nên ``im phăng phắc,'' trong khi con đường phía
bên kia sông thì chen chúc xe cộ, đông đúc đến mức hỗn loạn.

Theo truyền thống gia đình tôi, đêm 30 Tết cả nhà thường di chuyển sang nhà ngoại ở
huyện \ruby{Toei}{Đông Anh}\footnote{Tính đến 1/2025, huyện Đông Anh (Hà Nội) vẫn chưa được lên quận.} để ăn bữa cơm tất niên quây quần với nhà ngoại. Nhưng năm nay, với sự
hiện diện của gần như toàn bộ công ty \hll, nhà tôi phải linh hoạt hơn. Thay vì di
chuyển cả đoàn người, chúng tôi tận dụng ``Cánh cửa thần kỳ'' để đưa ông bà ngoại và gia
đình bên đó đến nhà tôi, cùng chuẩn bị một bữa tất niên đặc biệt.

Bữa tiệc tất niên năm nay, cũng giống như mọi năm, luôn được định hình quanh món lẩu,
nhưng lần này chúng tôi cũng sẽ làm các món ăn đặc trưng trong ngày Tết, và bánh chưng
là một trong số đó. Tuy nhiên, trong tay các cô gái \hll\ những món ăn này không
còn giữ nguyên hương vị (hoặc hình dáng) gốc. Nhưng thôi, đó là câu chuyện tôi sẽ kể
chi tiết hơn sau.

Vỉa hè trước nhà tôi trở thành sân khấu chính. Vì đường phố khá yên tĩnh, mọi người đã
trải chiếu (lấy từ nhà ngoại tôi - chắc lại một màn ``mượn tạm'') ra lòng đường để tụ tập.
Gia đình tôi thì khôn ngoan hơn, ngồi trên vỉa hè để đảm bảo an toàn.

Các dụng cụ làm bánh được mang đến đầy đủ, không thiếu thứ gì:

\begin{itemize}
    \item Nồi hấp bánh;
    \item Bộ chày giã cua;
    \item Hai xếp lá dong, cái nồi to nặng mà tôi đã ``xin'' được từ trong quê;
    \item Một ba lô củi gỗ do chính tay Noel-chan, Haato-chan và Rushia-chan hì hục chặt suốt cả buổi sáng;
    \item Nguyên liệu làm bánh như gạo nếp đã ngâm sáu tiếng, thịt lợn được ướp đậm đà từ hôm qua, và đậu xanh chín vàng ươm mà Miko-chi bằng cách nào đó có được.
\end{itemize}

Đích thân bố tôi và bà ngoại (mà từ giờ tôi sẽ gọi tắt là ngoại, tên bà ấy là Musada Aiko
(\ruby{無}{む}\ruby{佐}{さ}\ruby{田}{だ}\ruby{愛}{あい}\ruby{国}{こ}, \ruby{Mu}{Vô}-\ruby{xa}{Tá}-\ruby{đa}{Điền}\ \ruby{Ai}{Ái}-\ruby{cô}{Quốc})) đảm nhận vai trò ``sensei,'' hướng dẫn cả nhóm làm bánh
chưng đúng chuẩn. Với phong thái điềm đạm, họ chỉ bảo từng bước một, khiến không
khí vừa hào hứng vừa nghiêm túc.

Tuy nhiên, không phải ai cũng được tham gia. Từ phía xa, một giọng nói quen thuộc vang
lên từ trong nhà: ``Haachama này tại sao lại không được dự chứ!? Thật là bất công mà!''

Chắc các bạn cũng đoán ra ai rồi. Nếu các bạn còn nhớ việc tôi kể rằng em ấy đã từng
làm một chiếc bánh mì không thể ăn được\footnote{Xem {\flaregothic ÉPISODE 32}, quyển số Ba.}, thì đó vẫn chưa phải tệ nhất đâu. Tôi còn
nhớ em ấy còn từng làm cháy đen cả món sủi cảo được rán bằng dầu ô-liu kia mà. Tận
HAI lần.

Nhưng đừng lo, màn ``trừng phạt'' của chúng tôi với em ấy sẽ sớm khiến cả nhóm cười ra
nước mắt.

\ct

Với mọi thứ đã chuẩn bị sẵn, cả nhóm ngồi quây lại thành một vòng tròn trên chiếu. Bà
ngoại tôi bắt đầu hướng dẫn từng bước một, ánh mắt tràn đầy kiên nhẫn và nhiệt huyết.

``Trước tiên, các cháu sẽ xếp lá dong như thế này\dots'' bà bắt đầu lấy một xấp lá dong gồm
ba lá đặt lên đùi. Tôi tỏ vẻ tò mò: ``Bao nhiêu lá là đủ hả bà?''

Bà ngoại tôi đáp: ``Khoảng chừng hai, ba lớp là vừa.'' Tôi lấy đúng số lượng theo như bà
đã hướng dẫn.

Nhìn quanh, tôi thấy hầu hết mọi người làm khá ổn theo chỉ dẫn của bà. Ngoại trừ một
trường hợp nổi bật.

``Thế này đúng chưa ạ?'' giọng của Fubu-chan vang lên, với một xấp lá dong cao quá mặt
khiến ai nấy đều trố mắt. Tôi phản ứng ngay: ``Thế này thì nhiều quá! Em định làm món
lá dong chưng đấy á!?''

Tiếng cười vang lên khi Fubuki cười ngượng, vội vàng bỏ bớt.

\ct

``Xong rồi, các cháu lần lượt cho gạo nếp, thịt lợn, đậu xanh vào\dots'' đến lượt bố tôi tiếp
lời bà ngoại hướng dẫn cho mọi người.

Tôi nhìn sang Noel-PON, thấy em ấy cầm một tảng thịt to ngang lưng lên, làm cho chúng
tôi đổ mồ hôi hột. ``Cháu có thể cho cả tảng thịt to như thế này được không ạ?''

Ru-chan phản ứng: ``Tham quá đó, Noel'' Mel-chan thêm vào: ``Nếu thế thì sao mà gói
được đây?''

Noel-PON chỉ cười trừ, nhưng ánh mắt có vẻ vẫn còn luyến tiếc tảng thịt ``hoành tráng''
kia.

\ct

``Xong rồi các cháu gói lại thế này\dots\ Kuon ơi, lấy cho bà mấy cái dây chun ra đây.''

Mio-chan nhìn chiếc bánh được gói lại gọn gàng bằng những động tác phức tạp, nhíu mày
e ngại: ``Trông có vẻ khó nhỉ.''

Tôi gợi ý: ``Mọi người có thể dùng khuôn vuông để gói cho dễ mà.'' Tất nhiên là sau đó
lại có những người sẽ lại không theo quy tắc.

\begin{dialogue}
    \speak{Matsuri} Em có thể dùng không tam giác được không?
    \speak{Kuon} Không.
    \speak{Ayame} Thế hình quỷ thì sao nè?
    \speak{Kuon} Cũng không nốt. Cứ gói hình vuông đi là được rồi-
    \speak{Shion} Hình con chim origami!
\end{dialogue}

Phù thủy-chan giơ lên chiếc bánh chưng hình con chim hồng hạc được gấp như origami.
Tôi nhướng mày ngạc nhiên: ``Em gói nó lại kiểu gì vậy?''

Mẹ tôi xen vào, tỏ vẻ nghiêm nghị: ``Các cháu gói như thế thì cogn gì là bánh chưng nữa!''
Choco-sensei đắc ý kết luận: ``Như Kuon-mama-sama nói đó. Bánh chưng phải vuông
vắn mới đúng ý nghĩa.''

Mel-chan cũng gật gù đồng tình: ``Đúng vậy, đó là biểu tượng của đất mà. Đừng có phá
hỏng hay méo mó tinh thần của nó chứ.''

Trước lời khuyên của họ, hai nàng nhí nhố Shion-chan và Ayame-chan đành chịu thua,
đồng thanh đáp: ``Hai\~{}''

\ct

``Cuối cùng là dùng dây lạt buộc lại là xong,'' bố tôi lấy những chiếc dây lạt, buộc lại một
cách điêu luyện. Trông ông ấy có vẻ rất hào hứng, khi đã lâu rồi chúng tôi không được
gói bánh chưng như vậy.

Sora-chan chỉ vào những chiếc dây lạt, tỉ mỉ hỏi: ``Cháu nghĩ chúng ta buộc các đường
ngang và dọc xen kẽ nhau, như thế này đúng không chú?''

Ông gật đầu: ``Đúng rồi đó.''

Okayu-chan nhìn chiếc bánh chưng buộc xong, nhìn hình dáng của nó có vẻ quen thuộc:
``Nhìn giống bảng chơi cờ ca-rô ghê\dots'' Thấy thế, Korone-chan phấn khích: ``Ta thử chơi
trên đó đi!''

Cặp Chó-Mèo liền tận dụng nguyên liệu thừa để ``vẽ'' bảng caro trên một chiếc bánh và
bắt đầu chơi ngay tại chỗ. Kết quả, nàng Miêu tím dễ dàng đánh bại bạn mình chỉ sau bốn
nước.

``Tôi thắng rồi meo\~{}'' nàng Mèo thốt lên đầy thích thú. Nàng Khuyển trố mắt ra, tỏ vẻ
không phục: ``Ể\~{}? Ogayu, không công bằng! Chơi lại!''

Pekora ngồi cạnh đó, liền lên tiếng nhắc nhở: ``Đồ ăn không phải là đồ chơi đâu, peko!
Nghiêm túc giùm đi cái!''

Tất nhiên, còn có những trường hợp khó đỡ khác nữa. Ví dụ, Miko-chan thay vì gói bánh,
em ấy lại mải mê chụp ảnh tự sướng với chiếc bánh chưng của ``mình'' (thực chất ra là lấy
của Ru-chan) để đem khoe.

Hay, Aqua-chan thì vô tình gói luôn cả một cái thìa vào bên trong khi đang buộc dây,
khiến tất cả chúng tôi đều bật cười và cùng đồng thanh ``quê chữ ê kéo dài'' khi phát hiện.

Và, đoán xem người duy nhất không được tham gia vào công việc gói bánh chưng này
phản ứng gì nào? Haachama-chan thò đầu ra, hét lớn dù vẫn đang dỗi vì bị cấm cửa: ``\textbf{Các
    người sẽ hối hận vì không cho Haachama này tham gia!}''

Đáp lại chỉ là những tràng cười của tất cả chúng tôi, khiến cả một đoạn đường vốn im lìm
giờ đây tràn ngập niềm vui và sự náo nhiệt của gia đình.

\ct

Khi mọi người hoàn thành việc gói bánh, tiếng reo hò vang lên: ``\textbf{Đã xong!}''

Trước mặt họ, những chiếc bánh chưng đủ kiểu dáng, từ chuẩn chỉnh đến kỳ quặc, chất
thành một đống đầy màu sắc.

\uline{Từ những chiếc bánh vuông vắn, đúng chuẩn mực\dots}

Mẹ tôi cầm một chiếc bánh lên, trầm trồ khen ngợi: ``Miko-chan khéo tay thật đó!'' Vu
nữ-chan cười đắc ý: ``Chuyện nhỏ như con thỏ đó, ne!''

Tôi nhìn chiếc bánh của Sora-chan, tấm tắc khen: ``Cái của Sora-chan nhìn cũng chuẩn
phết đấy chứ.'' Em ấy nghe vậy, đáp lại bằng một nụ cười dịu dàng: ``Thế sao? Cảm ơn
anh nha!''

Bố tôi bổ sung thêm: ``Chú thấy Rushia-chan, Mel-chan với Mio-chan cũng khá đẹp đấy.
Tay nghề tốt thật.''

Mẹ tôi nhìn sang chiếc bánh mà cô em gái tôi đã gói: ``Ryuuhou cũng khéo tay ra phết
đấy chứ!'' Cô em gái tôi nghe vậy, cười lớn: ``Con mà lị!''

\dots Chẳng biết có phải ``con hát mẹ khen tay'' không, nhưng tôi phải thừa nhận là em ấy
gói gọn và đẹp thật.

\uline{\dots đến những chiếc méo mó khó hiểu\dots}

Bố tôi nhìn qua chiếc bánh mà tôi gói, thẳng thắn nhận xét: ``Kuon làm có vẻ hơi méo
nhỉ?'' Tôi chỉ có thể cười gượng trước mọi người: ``Mấy cái này con có rành đâu ạ\dots''

Thấy thế, Aki-chan giở giọng trêu: ``Ể? Em tưởng hôm nọ anh nặn được con thỏ\footnote{Phần {\abv Truyện phụ}, {\flaregothic ÉPISODE 35}, quyển số Ba.} mà?''

Tôi đáp nhẹ, khẽ thở dài: ``Hai cái đó khác nhau hoàn toàn, em ạ. Anh tốt cái này thì sẽ
phải dở phần khác chứ.'' Tôi liếc sang chiếc bánh của Noel-chan, nhận xét: ``Anh thấy
cái của Noel-PON còn phình to ra như đang mang chửa kia kìa!''

Nàng Hiệp sĩ cười khì, chẳng ngại ngùng mà khai nhận: ``Tại em thích ăn thịt, nên là\dots''
Ngoại tôi chỉ biết lắc đầu: ``Thì, đành vậy chứ sao giờ.''

Đến lượt mẹ tôi liếc sang chiếc bánh mà Akutan gói: ``Cái của Aqua-chan trông có vẻ bị
rách rồi kìa.'' Nghe vậy, em ấy hoảng hốt: ``Ơ!? Chết rồi, chắc là tại cháu gói vội quá\dots''

Tôi vội lên tiếng trấn an: ``Anh nghĩ em chỉ cần lấy một cái lá dong đắp vào rồi buộc thêm
dây là ổn mà. Em không phải lo đâu.''

\uline{\dots thậm chí còn có những chiếc ``tuyệt phẩm'' gây ngỡ ngàng, tới mức mà ``Quảng Nổ\footnotemark'' có thể thốt lên: ``Thật không thể tin nổi!''} \footnotetext{Tức Nguyễn Tử Quảng, CEO của BKAV.}

Tôi cầm lên chiếc bánh chưng của Ayame-chan, phát hiện nó nặng đến nỗi làm tay mình
ngoặt xuống. Ngạc nhiên, tôi hỏi: ``Em gói cái gì trong này mà nặng thế!?''

Ojou-san tỉnh bơ: ``Em gói giống như các anh hướng dẫn thôi, da zo.'' Tôi lại thắc mắc:
``Đúng hướng dẫn mà sao nó nặng vậy?''

Em ấy tiếp tục giải thích: ``Có cả kim loại để giữ khung mà.'' Tôi bật ngửa, đáp lại: ``Trời
ơi! Cần gì phải khổ sở thế!? Anh còn tưởng em cho cục đá vào đây đấy!''

Shion-chan nghe vậy, liền xen vào bình luận: ``Chắc là cái của bà để ném chết chó nhỉ?''
Tôi gật gù, phì cười với lời bình của em ấy: ``Cũng có thể lắm-''

``Mà khoan, từ từ!'' Quay sang nàng Phù thủy, tôi sững sờ khi thấy hai chiếc bánh chưng
của em ấy tỏa sáng rực rỡ: một chiếc phát ra ánh sáng trắng tinh khiết, chiếc còn lại u ám
như chứa đựng điều gì đó hắc ám.

Mẹ tôi cất giọng tò mò: ``Sao lại một cái lóa trắng, một cái lóa đen thế kia?'' Em áy đáp
với vẻ tự hào: ``Cái lóa trắng là \uline{bánh chưng của sự sống}, cái kia là\dots''

Tôi vội vàng ngắt lời: ``Để anh đoán\dots\ \uline{bánh chưng chết chóc}, đúng không?'' Shion-chan
reo lên: ``Pín-pon! Chuẩn luôn!''

Làm tôi nhớ đến một phân cảnh nào đó trong ``Chương trình thường nhật\footnote{Cụ thể hơn, là tập \textit{Sandwich of Death} (Bánh mì kẹp chết chóc), tập 4-13 của \textit{Regular Show} (Chương trình thường nhật.)}.'' Tôi thở dài:
``Shion-chan, làm ơn gói cái bánh đó một cách bình thường giùm anh đi\dots''

Mọi người bật cười sảng khoái, không thể ngừng bàn tán về những chiếc bánh chưng độc
nhất vô nhị của mình. Một vài người còn tranh thủ chụp ảnh lại để lưu giữ kỷ niệm đặc
biệt này.

\ct

Bước cuối cùng là chúng tôi sẽ đi đồ bánh. Tuy nhiên, chúng tôi lại gặp phải một vấn đề
không hề nhỏ: đó là số lượng bánh chưng mà chúng tôi gói lại quá nhiều, chưa kể đến
chiếc bánh siêu to khổng lồ (to gấp mười lần một chiếc bánh chưng thông thường) mà
Robo-PON gói nữa.

``Tuyệt lắm, Roboco-senpai ạ, giờ ta còn chẳng còn chỗ để bánh nữa\dots'' Mio-chan nhìn
gói bánh chưng khổng lồ mà em ấy làm, chỉ biết lắc đầu ngán ngẩm.

``Ít nhất thì em ấy còn biết gói bánh vuông vức,'' tôi nói, nhận xét một cách đầy công tâm.
Đáp lại, em ấy chỉ cười khì với chúng tôi: ``Xin lỗi mọi người nha, tại mình làm hăng
quá\dots''

Tôi quyết định để nàng Phù thủy trợ giúp. ``Shion-chan, em có thể hóa lớn cái nồi này
được không?''

Em ấy mỉm cười: ``Được chứ! Chuyện nhỏ! Xem nào\dots\ hô biến!'' Ngay lập tức, chiếc
nồi kim loại vốn dĩ đã lớn, giờ biến thành một chiếc nồi khổng lồ, đủ để chứa tất cả số
bánh chưng, kể cả chiếc bánh siêu to mà Roboco-chan tự hào gói.

Thực ra là cái nồi đó quá lớn, đến mức chiếm dụng luôn một phần đường đi rồi.


Roboco-chan thấy thế, cười tươi rói: ``Cảm ơn Shion-chan nhé! Bây giờ thì cái bánh của
chị cũng vừa rồi!''

Mọi người xúm lại, cẩn thận xếp từng chiếc bánh vào nồi. Cảnh tượng vừa náo nhiệt vừa
rộn rã tiếng cười đùa. Với chiếc bánh siêu khổng lồ của Robo-PON, đích thân em ấy và
Noel-chan phải khiêng vào cái nồi đó.

Nhưng khi đến công đoạn cuối cùng là đổ nước và nhóm lửa, không khí bắt đầu trầm
xuống khi họ nhận ra\dots

``Ta sẽ phải chờ khá lâu đấy,'' mẹ tôi cất giọng, tỏ vẻ hơi chán nản. Miko-chi nghe vậy, tỏ
vẻ tò mò: ``Bao lâu vậy cô?''

``Bình thường là sẽ khoảng 12 tiếng,'' bố tôi trầm ngâm một lúc lâu, rồi trả lời một cách
khá dứt khoát.

Cả nhóm - đặc biệt là Noel-chan, Roboco-chan, Akutan và Ru-chan - đồng thanh thốt lên:
``Lâu đến thế ư!?'' Trong khi đó, Ayame-chan chỉ nhún vai: ``Tôi nghĩ cái này giống như
ta đồ bánh mochi thôi mà, yo.''

Tôi chợt nhận ra một điều khá quan trọng, và lập tức hỏi bố tôi: ``Nhưng bố nói là \uline{bình
    thường}, còn thực tế mất bao lâu ạ?''

Ông ấy trả lời, lần này ngập ngừng: ``\dots Chắc là nguyên một ngày. Không cẩn thận là mất
ngày rưỡi.'' Câu trả lời này khiến mọi người bật ngửa: ``T-Thật luôn!?''

Shion-chan nhíu mày: ``Lâu thế cơ đấy\dots\ Thôi được rồi, để tôi lo việc này.'' Em ấy bước
đến gần nồi bánh chưng đang cháy sáng rực lửa, đưa tay thực hiện một câu thần chú:
``Phép thuật thời gian\dots\ Hô biến!''

Một luồng sáng bắn ra từ bàn tay Shion-chan, nhưng\dots\ chẳng có gì thay đổi cả. Cái nồi
bánh vẫn y nguyên như vậy, không có một hiệu ứng đặc biệt nào hiện ra cả.

Tôi nhướng mày, đoán được ý đồ của em ấy: ``Em đang định tua nhanh thời gian cái nồi
này sao?'' Shion-chan cười gian: ``Ái chà, em không nghĩ là anh thông minh đến thế đấy.''

``Nghe cách em hô biến như thế thì ai mà chả đoán được, ne,'' Miko-chi nói với vẻ mặt
không hề biến sắc gì. Tất cả mọi người cũng gật đầu đồng tình với em ấy.

Tôi nhìn xung quanh cái nồi đó, nhưng vẫn không nhận ra được dấu hiệu gì cho thấy nó
đang được tua nhanh. ``Sao anh chẳng thấy gì vậy?''

Shion-chan chỉ vào đống lửa bên dưới nồi: ``Cái lửa đang bập bùng cháy nhanh hơn thế
thôi.'' Tôi ghé sát xuống đống củi, và quả nhiên, ngọn lửa trông như đang cháy với tốc độ
gấp đôi bình thường. Chưa kể rằng khói bốc lên còn nhanh hơn nữa.

Nhưng tôi vẫn cảm thấy bán tín bán nghi, bèn hỏi tiếp: ``Em tua nhanh đến đâu vậy?''

``Gấp đôi.'' Fubuki-chan lên tiếng: ``Sao em không tăng gấp mười lần nhỉ?''

Lập tức, nàng Phù thủy phản ứng: ``Em cạn kiệt ma lực rồi. Cũng tại chị nhờ em dọn cái
tầng ba chứ bộ. Gấp đôi là nhanh rồi đấy.'' Nàng Cáo trắng chỉ biết cười trừ: ``Quả thật bọn chị có nhờ em hơi nhiều nhỉ\dots''

Câu trả lời không mấy tích cực này khiến cho Choco-san thở dài: ``Vậy là ta vẫn phải chờ
ít nhất là 12 tiếng để bánh chín à\dots'' Mel-chan cười méo xệch: ``Chưa kể phải luôn canh
bếp nữa chứ\dots''

``Em bắt đầu thấy đói rồi\dots'' Koro-chan than vãn, vờ ôm bụng quằn quại. Okayu-chan
cũng từ đó mà hùa vào: ``Em cũng thế\~{}''

Tôi gật gù thừa nhận: ``Ừ. Gần như ai cũng đói mà. Nhưng mà ta vẫn cần một ai đó để
canh cái nồi bánh chưng này.''

Tôi nhìn quanh tìm kiếm ai đó để giao việc, và chợt nhớ ra người cuối cùng chưa tham
gia gói bánh: ``Akai Haato-chan.''

Tôi gọi lớn, đủ để vọng qua cửa: ``\textbf{Haachama-chan! Bọn anh có việc nhờ em này!}''
Ngay tức khác, nàng Trung học chạy ra ngoài, với vẻ phấn khích tột độ: ``Cuối cùng cũng
đến lượt Haachama ra tay rồi à!?''

Tôi cười mỉm: ``Ừ. Anh có việc nhờ em đây.'' Haato-chan càng sáng mắt hơn: ``Việc gì
cũng được tất!''

\ct

``\dots Tại sao mình phải làm việc này chứ?'' Nghe thấy tiếng than thở vọng lên, tôi lại ngó
đầu ra ngoài nhìn vỉa hè trước cửa nhà.

Đây đã là lần thứ en-nờ Haato-chan đã than thở như vậy. Cứ mỗi lúc, tiếng than trách bản
thân càng lúc càng lớn hơn, đến độ vang vọng khắp trời.

Em ấy đang ngồi ngoài sân, một mình canh nồi bánh chưng khổng lồ. Công việc của
Haato-chan rất đơn giản (nhưng cũng chán ngắt): cứ mỗi 15 phút lại thêm củi vào lò, đảm
bảo ngọn lửa cháy đều.

``Chán quá đi mất thôi\dots\ Ngoài này còn lạnh nữa chứ\dots'' em ấy lại than vãn, thi thoảng
còn hắt xì mấy cái rõ to.

Giữa trời tối, với một ngày mây khá dày đặc, chỉ còn mỗi nàng Song tính vừa cặm cụi đốt
củi, vừa tự lẩm bẩm: ``\textit{Tại sao tôi lại không được gói bánh chưng với mọi người chứ!?}''
Em ấy chỉ còn biết ngồi gần tới nồi bánh chưng đang bốc lửa để giữ ấm bản thân, trong
khi phải chống chịu với cái đói và cái rét đậm của miền Bắc.

\ct

Tám giờ tối.

Bữa tất niên của gia đình chúng tôi đã kết thúc trong không khí ấm cúng và vui vẻ. Đây
là một bữa ăn hoàn toàn bình thường - không phá hoại, không trò đùa quá đà, và cũng
chẳng có pha xử lý mất não nào xuất hiện.

Tuy nhiên, giữa không khí rôm rả ấy, vẫn có một người bị bỏ lại.

Đó là cô nàng tóc vàng ruy băng đỏ - Akai Haato - ngồi trông nồi bánh chưng với vẻ mặt
chán nản. Thỉnh thoảng, em ấy đứng dậy, đi lại mấy vòng quanh nồi lửa để giết thời gian.

``Mọi người xong chưa vậy không biết\dots'' tôi nghe em ấy thở dài, mắt lơ đãng nhìn lên
bầu trời đêm.

Tôi nghe thấy tiếng thoang thoảng bụng của em réo lên biểu tình, đi kèm với đó là một
câu than vãn nữa: ``Đói quá\dots\ Dù ở đây có lửa ấm nhưng sao vẫn thấy lạnh trong tim
thế\dots''

Em ấy đưa tay nắm lấy ngực, khuỵu xuống đất với vẻ mặt bi kịch: ``Tạm biệt mọi người,
mọi người đúng là quá đáng\dots''

Tôi không nhịn được cười nữa. Cuối cùng, tôi bước ra khỏi cửa, cất giọng hỏi: ``Làm cái
gì mà đóng kịch kinh thế?''

Nghe giọng của tôi, Haato-chan giật mình đứng phắt dậy, suýt đạp vào cái nồi bánh chưng.
``C-C-C-Cái gì thế!? Xém nữa là đổ nồi rồi\dots''

Mặt của em ấy đỏ hỏn lên, dường như vì cảm thấy ngượng ngùng khi ai đó nhìn em ấy
phản ứng như vậy.

Tôi đưa chiếc cặp lồng đựng thức ăn còn nóng hổi, rồi mỉm cười trừ: ``Xin lỗi. Bọn anh
mang đồ ăn xuống cho em đây.''

Mắt Haato-chan sáng lên khi thấy cặp lồng lẩu bốc khói nghi ngút. Em lí nhí cảm ơn:
``Cảm ơn anh\dots'' nhưng rồi bất giác nhận ra điều gì đó, em ngước lên hỏi: ``Khoan, bọn
anh?''

Chưa kịp dứt lời, một loạt giọng nói vang lên phía sau. ``Bọn này cũng muốn xuống nữa,
yo!'' Ayame-chan nhón bước tới gần, cất lên một giọng nói tinh nghịch.

Fubu-chan đi ngay phía sau, tủm tỉm cười: ``Một mình trông nồi bánh chắc là cô đơn lắm
nhỉ?'' Subaru-chan thêm vào, tay khoanh trước ngực: ``Nhìn chị ngồi một mình thế này
là biết rồi.''

Nhìn chúng tôi cười khẩy bày tỏ sự quan tâm như thế, Haato-chan lúng túng, mặt đỏ bừng,
rồi buông ra một câu mang đúng tính ``tsundere'' của mình: ``A-Ai khiến mọi người xuống
đây chứ!''

Mel-chan nheo mắt, cất giọng bông đùa: ``Thế chúng ta lại lên vậy.'' Tôi cố hết sức để
nhịn cười, đáp: ``Anh cũng đang định nói thế!''

Vừa dứt câu, mọi người giả vờ quay người bước lên nhà. Hoảng hốt, nàng Trung học giơ
tay ngăn lại: ``K-Không! Ý của tôi không phải thế!'' và rồi hét to, đầy tuyệt vọng: ``\textbf{Mọi
    người quay lại đây đi mà!}''

Chúng tôi phá lên cười. Nhìn vẻ mặt đẫm lệ của Haato-chan vì bị chúng tôi trêu, Quỷ
sừng nhẹ nhàng đặt tay lên vai em ấy, nói: ``Thôi, đùa thế thôi. Bọn em cũng sẽ xuống
đây ngồi canh bánh chưng với chị đây.''

Cáo trắng-chan lôi ra vài chiếc ghế nhựa mà em ấy đã chuẩn bị sẵn: ``Bọn tôi cũng sẽ làm đồng
minh canh lửa với bà đêm nay! Không cần phải lo đâu!''

``Đúng vậy, bọn em sẽ không để chị một mình đâu,'' Subaru nói, tỏ vẻ nghiêm túc với
quyết định của mình.

Haachama-chan thở phào, ngồi phịch xuống ghế, long lanh nhìn chúng tôi: ``Mọi người\dots\ cảm ơn nhé.''

\ct

Tại Etsunan thời xưa, việc trông nồi bánh chưng không chỉ là một nhiệm vụ mà còn là
một thú vui, nơi cả gia đình quây quần bên nhau, tham gia vào các hoạt động ấm áp, đậm
chất truyền thống.

Ở thời hiện đại, tôi vẫn giữ nguyên tinh thần đó, đặc biệt là trong bối cảnh ngày càng ít
gia đình còn nấu bánh chưng như xưa. Nhưng với nhóm chúng tôi, không khí xung quanh
nồi bánh lại có phần khác biệt, và có phần khó tả.

Một PON - đúng vậy, Noel-chan - đang tập thể dục bằng chiếc chùy sắt mà em cứ giữ khư
khư trên người. Một đại PON khác, Roboco-chan, đang tự tháo cánh tay mình ra để bảo
dưỡng, trông chẳng khác gì một buổi ``mổ xẻ robot'' tại chỗ.

Ở một góc khác, Miko-chan đang hăng say biểu diễn ``điệu múa thần linh'' tự biên tự diễn,
với bộ kimono đỏ rực cuốn hút ánh mắt mọi người. Cặp đôi Chó-Mèo - Korone-chan và
Okayu-chan - thì chơi trò đuổi bắt, kéo theo Ayame-chan và Shion-chan nhập hội, biến
khung cảnh quanh nồi lửa thành một màn rượt đuổi náo loạn.

Mel-chan, Subaru-chan, và Haato-chan thì lập thành một nhóm chơi bài, trong khi Miochan
và Fubuki-chan chọn cờ lật làm trò tiêu khiển dù trời đã tối đen. Các thành viên
còn lại, bao gồm tôi và vài người nữa, ngồi thành từng nhóm nhỏ, buôn chuyện linh tinh,
chẳng khác nào các nam sinh nữ tú trung học tụ tập sau giờ học.

Riêng A-san, như thường lệ, cầm máy quay đi khắp nơi ghi lại từng khoảnh khắc hài hước.

Gia đình tôi thì yên tĩnh hơn đôi chút. Cả nhà ngồi quanh nồi bánh, mắt dán vào màn hình
TV, cùng thưởng thức chương trình mà bất cứ ngườiđiện thoại miền Bắc nào cũng mong
chờ mỗi dịp Tết đến xuân về - Táo Quân, hay \textit{Gặp nhau cuối năm}.

Pekora-chan ngồi gần đó, nghiêng tai thỏ của mình nghe tiếng cười của chúng tôi, rồi tò
mò hỏi: ``Có cái gì mà mọi người cười kinh thế, peko?''

Tôi quay sang đáp: ``Bọn anh đang xem Táo Quân ấy mà.'' Thấy Thỏ-chan ngơ ra và định
hỏi thêm, Sora-chan và Mel-chan liền giải thích:

\begin{dialogue}
    \speak{Mel} Nó dựa theo truyền thuyết Táo quân chầu trời, Pekora-chan.
    \speak{Sora} Đúng đó. Theo chị được biết, vào ngày này, tất cả các Táo lên trời dự buổi
    chầu, báo cáo Ngọc Hoàng các sự kiện sau một năm.
\end{dialogue}

Peko-chan lại chớp mắt: ``Làm sao hai người biết vậy, peko?'' Sora-chan mỉm cười: ``Cả
tối hôm qua bọn chị lên mạng tra cứu mà.''

Lúc này, Aqua ló đầu vào, hào hứng tiếp lời: ``Em cứ tưởng tượng giống như Fubuoshi
chúng ta thôi\footnote{Xem {\flaregothic ÉPISODE 9}, quyển số Một.}!'' Thấy ai đó nhắc đế mình, Fubuki ngoái đầu lại, gương mặt đầy thắc
mắc: ``Ai đó nhắc chị à?''

Nàng Thỏ quay sang tôi, mắt đầy tò mò: ``Cơ mà nếu là báo cáo thì chẳng phải anh ngày
nào cũng làm thế sao, peko?''

Tôi bật cười, giải thích thêm: ``Đấy là báo cáo giấy. Còn trong chương trình này, người
ta báo cáo kết hợp cả hài. Kiểu như ngâm thơ hay là hát chẳng hạn.''

Thỏ-chan nghe vậy, càng cảm thấy khó hiểu hơn, tưởng tượng ra cảnh tôi ngồi viết báo
cáo bằng thơ hoặc hát để chọc cười A-san và Tanigo-san. Em liền hỏi: ``Như thế sao,
peko?''

Tôi gượng cười, đổ mồ hôi: ``Có thể hiểu là như vậy. Mà công ty nhà ta lắm chuyện ngớ
ngẩn đến mức có thể viết báo cáo theo tuần chứ chẳng cần theo năm đâu.''

Akutan nghe vậy, e thẹn gật đầu: ``Anh nói đúng quá\dots''

``Thậm chí là bây giờ chẳng hạn,'' tôi khẽ nói. Nàng Dơi nghe vậy, liền lên tiếng thắc
mắc: ``Vậy anh thử nói xem, tình huống nào có thể khiến báo cáo hài xảy ra ngay lúc này
được?''

Tôi không nói gì, chỉ tủm tỉm cười, tay chỉ về phía mọi người xung quanh để trả lời câu
hỏi của em ấy.

Ở một góc đối diện đường, Noel-chan đang cố gắng giúp Roboco-chan sửa chữa theo cách
không giống ai: bằng cách dùng chiếc chùy mà nàng Hiệp sĩ vừa cầm nãy giờ đập mạnh
vào lưng em ấy.

Thấy nàng Hiệp sĩ cứ liên hồi đập, Robo-PON thốt lên: ``Ái da, đau! Em đang làm cái gì
thế!?''

Noel-chan cười mỉm: ``Em đang sửa chị đó, Roboco-san! Chị cứ ở yên để em đập thêm
phát nữa\dots'' Roboco thấy vậy, hoảng loạn thốt lên: ``Đập thế này thì chị hỏng luôn rồi còn
gì! Ái da!! \textbf{Không!!!}''

Ở góc khác, nhóm chơi đuổi bắt lúc đầu đã chuyển thành\dots\ màn chạy trốn sát nhân khi
Ru-chan, với con dao trên tay, hùng hổ đuổi theo Korone-chan và Okayu-chan.

Rushia-chan vừa điên cuồng vung dao, vừa hét lớn như đang hát nhạc metal: ``\textbf{Tại sao
    lúc nào tôi cũng là người thua chứ!? Mấy người đứng lại cho tôi ngay!!}''

Trong khi đó, cặp đôi Chó-Mèo thì không ngừng trêu tức hậu bối của mình, giống như
cách họ đang chọc một con chó dữ:

\begin{dialogue}
    \speak{Korone} Tại em chơi gà thôi\~{}
    \speak{Okayu} Đúng đúng, học hỏi bọn này đi chứ.
\end{dialogue}

Càng lúc, nàng Pháp sư càng tức điên hơn. Cô la lớn, tới mức lạc cả giọng: ``\textbf{Tôi sẽ giết
    chết mấy người!!! ĐỨNG LẠI!!}''

Thấy vậy, cặp đôi lập tức giở giọng châm chọc trêu tức hơn: ``Ta cùng chạy nhanh thôi\~{}''
Theo sau Ru-chan là Ayame-chan, người cũng đang hớn hở: ``Này, cho tôi chạy chung
với\~{}!''

Ở phía nhóm chơi bài ngay cạnh chúng tôi, hình phạt dành cho người thua đang diễn ra.
Tôi thấy Shion-chan đang đứng suy nghĩ, rồi bất chợt nói: ``Lần này là\dots\ cá voi!''

Subaru-chan tỏ vẻ ngỡ ngàng: ``Làm sao mà tôi biết con cá voi kêu như thế nào!?'' Haatochan
nói với sự xóc xỉa: ``Chị tưởng nó kêu 'ù\~{}' chứ nhỉ?''

``Nó đâu có kêu như thế!!'' Vịt-chan phản ứng lại. Shion-chan sau đó đã đưa ra một đề
bài dễ hơn: ``Thôi được rồi. Bây giờ bà kêu tiếng vịt đi.''

Dứt lời, cả Haato-chan và em ấy đều cố hết sức nhịn cười. Người đang bị chịu hình phạt
thấy em ấy đang nói đểu mình, liền vùng văng hét lên: ``\textbf{Đừng tưởng tôi sẽ kêu nha!}''

``Kêu gì?'' nàng Song tính hỏi, nở nụ cười trên khóe môi. Subaru-chan thở dài, rồi bắt
chước tiếng vịt kêu: ``Quạc, quạc, quạc.''

Ngay khi Subaru-chan vừa ``kêu tiếng bản xứ'' của mình, cả nàng Phù thủy và nàng Song
tính đều phá lên cười: ``Đấy! Bà/Em tự kêu đấy nha!''

Nàng Vịt giật mình khi nhận ra em ấy đã bị dính bẫy, vội thanh minh với hai người:
``Không! Không phải thế!!'' rồi cầm cái loa của mình đuổi theo họ.

\begin{dialogue}
    \speak{Haato} Ôi không\~{} Cô Vịt giận rồi kìa\~{}
    \speak{Shion} Chạy đi, chạy đi nào!
    \speak{Subaru} Tôi bảo rồi, tôi không phải là vịt!! \textbf{Đứng lại đó cho tôi!!}
\end{dialogue}

Tôi hy vọng bốn người họ không làm quá lên\dots\ Thậm chí tệ hơn còn bị lên đồn chứ chẳng
đùa.

Ở phía đối diện bọn họ, nhóm chơi cờ lật cũng đang có chuyện dở khóc dở cười không
kém.

Dù họ không cất lên một lời nói, nhưng Fubuki-chan và Mio-chan lại đang đổ mồ hôi hột
như suối, cả thân người run lẩy bẩy, gương mặt thì cố ra vẻ bình tĩnh, thỉnh thoảng còn
rên lên khó chịu.

Sử dụng khả năng ``đọc suy nghĩ động vật'' mà tôi học lỏm được từ anh ``Tô'', tôi lập tức
hiểu được lý do hai người họ lại im thin thít như vậy.

Mio-chan nghĩ trong sự hối thúc: ``\textit{Bà ngồi đó nửa tiếng rồi đấy\dots\ Đi nhanh lên, tôi sắp
    nhịn không nổi nữa rồi!}''

Fubuki-chan cũng không khá khẩm hơn: ``\textit{Đi đường nào cũng chết\dots\ Mà mình mót ị
    quá\dots}'' Liếc lại sang nàng Sói đen, em ấy thậm chí còn ngồi vắt chéo, thi thoảng còn giữ
chặt hạ bộ của mình: \textit{``Ôi mót quá\dots\ Nhanh lên, tôi sắp tè ra quần rồi!!}''

\dots Quả đúng là một cảnh tượng rất khó xử. Chả trách mà họ không dám ho he lên tiếng.

``Pekora-chan?'' tôi gọi nàng Thỏ vẫn đang bắt chuyện với Choco-sensei.

``Sao thế, peko?''

``Em với Choco-san có thể giúp họ đánh nốt ván cờ lật được không?''

``Là sao vậy, Kuon-sama?'' đến lượt Sensei hỏi. Tôi chỉ tay vào cặp đôi FubuMio đang
run rẩy đằng kia, rồi nói khẽ: ``\hl{Hai người họ muốn đi vệ sinh.}''

Nàng Thỏ và nàng Y tá hiểu được ý của tôi, đi đến gọi họ. Họ sau đó hoàn thành nốt ván
cờ, giải thoát cho hai người đang ``khổ sở với nỗi buồn.''

Trong khi đó, các thành viên còn lại vẫn tiếp tục buôn dưa lê như những nữ sinh trung
học. A-san vẫn hăng say quay phim, ghi lại từng khoảnh khắc hài hước. Không có gì bất
thường xảy ra với chúng tôi cả.

Sau khoảng hai tiếng ngồi canh bánh, mọi người quyết định tụ họp lại, biến nồi bánh
chưng thành tâm điểm của một buổi karaoke cuối năm. Bao nhiêu mối bất hòa giữa mọi
người sau đó cũng được giải quyết hết, không còn ai có ý kiến gì bực bội trong người nữa.

Bộ dàn âm thanh xuất hiện từ\dots\ hư không, và từng giọng ca vang lên giữa không khí ấm
cúng, náo nhiệt.

Một buổi tối cuối năm đầy sự tréo ngoe nhưng cũng không kém phần vui vẻ, để lại những
tiếng cười giòn giã bên nồi bánh chưng đang dần chín tới.

\ct

Mười một giờ năm mươi chín đêm.

Thời khắc chuyển giao giữa năm Ất Hợi và năm Canh Tý chỉ còn được tính bằng giây.
Tôi đứng đó, đầu óc lơ mơ như thể có điều gì rất quan trọng mà mình đã quên mất.

Thay vì để Sora-chan, người luôn dẫn đầu trong những dịp đặc biệt - hô gọi, Matsuri đã
hăng hái mà cướp lời, hào hứng nói với mọi người: ``Mọi người ơi, chuẩn bị đếm ngược
nha!''

Cả căn phòng như bừng lên khi tất cả cùng đồng thanh đếm ngược, chỉ có tôi là vẫn mãi
trầm ngâm cố gắng nhớ lại.

``\textbf{Mười! Chín!}''

Là cái gì nhỉ?

``\textbf{Tám! Bảy!}''

Cái gì thực sự quan trọng lắm\dots

``\textbf{Sáu! Năm! Bốn!}''

\dots Thôi kệ đi. Tôi cũng hòa mình vào bầu không khí rộn ràng và cùng nhau đếm ngược
với mọi người.


``\textbf{Ba! Hai! Một\dots}''

Đồng hồ chỉ đúng mười hai giờ đêm. Tiếng hô vang ``\textbf{Chúc mừng năm mới!}'' đồng thanh
vang vọng khắp phố phường, phá tan mọi cảm giác lặng lẽ của năm cũ.

Ngay sau câu chúc mừng, một loạt tiếng pháo hoa vang lên từ xa, vọng lại qua các tòa
nhà. Tất cả dừng lại, mắt hướng lên trời cao.

Nhưng\dots\ chúng tôi lại chẳng thấy cái gì cả, ngoài việc tôi nghe được những tiếng ``bụp,
bụp'' của pháo hoa, và một chút ánh sáng sắc màu được lóe lên trên trời cao.

``Tiếng pháo hoa thì nghe rõ đấy, nhưng mà\dots'' Roboco-chan lên tiếng, Choco-sensei cũng
cảm thấy kỳ lạ: ``Không thấy đâu cả nhỉ?''

Miko-chi giơ tay, chỉ về phía những ánh sáng xanh đỏ lập lòe phía sau một tòa cao tầng:
``Kia kìa, mọi người!''

Haato tỏ vẻ bất mãn: ``Nhưng nếu thế nhìn được cái gì chứ!? Chúng ta đang bị che mắt
hết rồi kìa!''

Đúng lúc này, tôi chợt nhớ ra điều đã làm mình bận tâm từ nãy đến giờ: ``Anh quên mất,
chúng ta phải lên mái nhà mới xem được pháo hoa.''

Tất cả mọi người sững sờ, rồi cùng nhìn về phía tôi với những ánh mắt trách móc. Matsuri-chan
thậm chí còn nổi cáu: ``Vậy tại sao anh không nói trước!?''

``Thì bây giờ anh mới nhớ ra mà\dots'' tôi ôm đầu nói, đồng thời hướng về mọi người và gửi
một lời xin lỗi.

Tôi vội chạy lên tầng sáu, kiểm tra xem mái nhà đã được mở chưa. Chiếc cầu thang gỗ
dẫn lên đó thì vẫn còn, nhưng cánh cửa gỗ thì bị khóa rồi.

Tôi lao đầu chạy xuống nhà, báo với bố mẹ: ``Bố ơi! Bác Hijaki-san năm nay không mở
căn mái ạ?'' Bố tôi thở dài, tỏ vẻ hơi tiếc nuối: ``Mấy năm nay bác ấy ốm yếu, nên không
mở thang lên mái nhà được.''

Mẹ tôi cũng gật đầu, tỏ ra bất ngờ: ``Giờ anh nhắc em mới nhớ. Hôm nay bác ấy còn bận
trông ông bà nữa.''

Tôi thở dài, nhìn qua các cô gái đang tỏ vẻ hụt hẫng: ``Chắc chúng ta chỉ có thể xem qua
ti-vi thôi nhỉ\dots''

Trong lúc chúng tôi tỏ vẻ chán nản vì có lẽ sẽ không thể chiêm ngưỡng được màn bắn
pháo hoa thường niên, Shion bất chợt lên tiếng: ``Sao không để em đưa tất cả mọi người
lên nhỉ?''

Tôi ngoảnh mặt lại, khẽ ngạc nhiên: ``Hả?'' Subaru-chan thậm chí còn không tin: ``Đừng
có ngớ ngẩn nữa, Shion. Ta cùng lên nhà thôi-''

Thế nhưng, chưa kịp để em ấy dứt lời, Shion đã bắt đầu tạo phép, đọc câu thần chú lặp
đi lặp lại nhiều lần. Theo sau đó, một vòng tròn màu xanh lam khổng lồ xuất hiện dưới
ngay chân chúng tôi, đủ lớn để bao quát tất cả mọi người.

``\textbf{Cuồng phong!}''

Một cơn gió giông chợt nổi lên, đánh bật chúng tôi lên trên cao. Bản thân tôi còn cảm
thấy ngỡ ngàng với cô nàng lắm chiêu này, còn các cô gái thì hốt hoảng hét lớn: ``\textbf{Chuyện
    gì đang xảy ra thế này\dots!!}''

Chẳng bao lâu, chúng tôi đáp xuống mái nhà của chúng tôi. Cũng may cho chúng tôi,
Shion-chan khá tinh ý khi nhẹ nhàng đưa chúng tôi chạm xuống mái nhà. Tuy nhiên, một
vài người trong số đó vẫn bị mất căn bằng và ngã bệt xuống đất.

\begin{dialogue}
    \speak{Pekora} Ái da da, peko\dots
    \speak{Subaru} C-Cô đang làm cái trò gì đấy hả!?
\end{dialogue}

Trong lúc mọi người còn không hiểu Phù thủy-chan đã làm gì với cúng tôi, Matsuri phấn
khích, nhanh chóng kéo mọi người dậy: ``Mau lên, mau lên nào! Mọi người ơi, trông đẹp
lắm á!''

Tất cả chúng tôi hướng về phía em ấy đang chỉ trên trời. Tại đây, bầu trời đêm ở Kawachi
- thủ đô của Etsunan hiện ra sống động với đủ sắc màu từ những chùm pháo hoa bắn lên
cao. Tiếng nổ lép bép và âm thanh reo hò hòa quyện vào nhau.

Aki-chan mỉm cười, đôi mắt long lanh nhìn lên bầu trời: ``Đây là lần đầu tiên chị được
xem màn pháo hoa ngoài Tokyo như thế này\dots\ Đẹp thật á\dots''

Lễ hội-chan thì sáng mắt lên: ``Đẹp quá\dots\ Quả này còn hoành tráng hơn cả lễ hội mùa hè
nữa.'' Akutan cũng công nhận: ``Công nhận họ chuẩn bị chu đáo ghê.''

Tôi ngước mắt nhìn những quả pháo hoa rực rỡ được bắn lên, rồi cùng gật gù đồng ý:
``Phải nói là năm nay thành phố chuẩn bị rất kỹ lưỡng đấy.''

Subaru-chan cũng đồng tình với tôi: ``Mọi người ở thành phố này hẳn 'đã đầu tư không ít
đâu. Đúng không, Haato-senpai?'

Nàng Trung học cười lớn: ``Ừm. Có ki còn ngang hàng với những buổi lễ ở thành phố lớn
đấy chứ.''

Toàn thể hai mươi tưtư người, mỗi người một phản ứng, đều cảm thấy rạo rực trong lòng
khi chứng kiến màn trình diễn pháo hoa đầy mãn nhãn này.

Nó đồng thời cũng thôi thúc trong lòng chúng tôi rằng một mùa xuân nữa lại tới. Dẫu
cho trời vẫn còn rét run cái lạnh đặc trưng của miền Bắc, tôi và mọi người đều có thể cảm
nhận được sự chuyển giao giữa hai mùa, ngay tại chính thời điểm này.

Chúc mừng năm mới, \hll!

\ct

Màn pháo hoa kéo dài chừng mười phút rồi dần kết thúc, để lại chút dư vị nuối tiếc trong
lòng mọi người.

Fubu-chan đưa tay lên trời, tỏ vẻ tiếc nuối: ``Ế? Hết rồi à? Đang vui thế cơ mà\~{}\dots''
Ru-chan thì phụng phịu: ``Tưởng phải lâu hơn chứ! Mới tí mà\dots''

Tôi nhìn vẻ thất vọng của mọi người, vừa cười vừa lắc đầu: ``Mấy đứa kỳ vọng nhiều quá
rồi. Chẳng phải pháo hoa thường chỉ kéo dài tầm đó thôi sao?''

A-san chỉnh lại cặp kính, đồng ý với bình luận của tôi: ``Ở Tokyo người ta cũng chỉ bắn
tầm ba tới năm phút. Thành phố Kawachi như thế này khá là rộng lượng rồi đấy.''

Nghe vậy, tất cả mọi người ngước nhìn lên bầu trời đêm, giờ đây đã được trả về vẻ yên
bình vốn có, khẽ mong ngóng một cái gì đó nhộn nhịp hơn.

Bỗng nhiên, Fubuki giơ tay như thể vừa nghĩ ra điều gì tuyệt vời: ``Này, tại sao chúng ta
không tự tổ chức một buổi pháo hoa của chính mình nhỉ?''

Mọi người đều đồng loạt quay sang nhìn em ấy, vẻ mặt từ bất ngờ chuyển sang phấn khích,
nhiệt liệt hưởng ứng một cách mãnh liệt.

Miko-chan thốt lên: ``Ý kiến hay đấy, Fubuki! Làm ngay thôi nào!'' trong khi Matsurichan
cũng tỏ vẻ mừng rỡ: ``Đúng á! Chứ chỉ có mười phút thế này, chán òm lắm!''

Chỉ riêng tôi và A-san thì đều nhướng mày, xen lẫn chút tò mò.

\begin{dialogue}
    \speak{A-chan} Vấn đề là các cô sẽ thực hiện nó kiểu gì?
    \speak{Kuon} Mà, mấy đứa định tự tổ chức pháo hoa thật à?
\end{dialogue}

Mẹ tôi nhẹ nhàng trấn an tôi: ``Cứ để họ làm xem nào. Có khi vui ấy chứ.''

Bố tôi mỉm cười, cố gắng giúp chúng tôi thoải mái hơn: ``Yên tâm đi, chắc họ cũng chỉ
chuẩn bị được vài quả pháo đất thôi. Không đến mức nổ banh mái nhà đâu.''

Và thế là, cả đám lao nhao bàn luận, chia nhóm, quyết tâm tổ chức một màn pháo hoa
``cây nhà lá vườn''. Còn tôi thì chỉ biết đứng nhìn, trong lòng vừa ngờ vực, vừa hào hứng
chờ xem màn trình diễn kỳ lạ này sẽ đi tới đâu.

\dots Chẳng ngờ rằng, đó sẽ là một trong những pha ngớ ngẩn và hỗn loạn nhất ngay từ khi
bước sang những giây phút đầu tiên của năm Canh Tý.

Fubuki-chan, với tinh thần ``hay chơi lớn'', en ấy mang ra một quả pháo hoa khổng lồ có
đường kính bằng một chiếc bàn học, kèm theo hàng chục quả lớn nhỏ chất đống xung
quanh.

Mio-chan nhìn cảnh tượng này, không giấu nổi vẻ hoảng hốt: ``Bà lấy mấy thứ này ở đâu
vậy!?'' Bạn thân của em ấy cười tươi rói: ``Ở kho vũ khí bí mật của Pekora-chan đấy!''

Peko-chan, đang đứng cách đó không xa, nghe thấy vậy liền giật mình: ``Này! Đống đó
là của tôi mà! Ai cho chị động vào hả, peko!?''

Nhưng chưa kịp giằng lại, Matsuri-chan, với sự năng nổ thường trực, đã nhanh tay châm
lửa vào dây pháo mà chẳng thèm nghĩ ngợi. Nàng Thỏ nhảy bổ tới, cố gắng thổi tắt dây cháy, nhưng mọi nỗ lực đều thất bại.

``\textbf{Nguy rồi, peko! Nguy rồi, peko!!}'' em hét lớn, hoảng loạn. Thế rồi, trong một pha xử
lý ``đi vào lòng đất'', Pekora-chan quyết định ôm chặt quả pháo khổng lồ như thể muốn
hy sinh để cứu lấy mọi người.

Sự việc diễn ra quá nhanh khiến mọi người không kịp căn ngăn, để rồi khi dây cháy hết,
quả pháo đó bay lên không trung với tốc độ cao.

``\textbf{Pekora(-chan)}!!!'' tất cả chúng tôi la lên thất thanh đầy lo lắng khi thấy quả pháo bùng
cháy, bay vút lên không trung mang theo nàng Thỏ xanh xấu số.

``Cứu tôi với!! Peko\~{}!!'' tiếng của em ấy vọng xuống khi quả pháo bắt đầu tỏa sáng rực
rỡ giữa bầu trời đêm.

\textbf{*BÙM!*}

Quả pháo nổ tung trên cao, tạo ra những dải sáng lấp lánh, theo sau là hàng loạt chùm
pháo hoa nhỏ bắn liên tiếp. Bầu trời đêm trở thành một bức tranh đầy màu sắc.

Lễ hội-chan, mắt sáng rực, không giấu nổi sự hào hứng: ``Lần này trông còn ăn đứt màn
pháo hoa vừa nãy!'' Ru-chan gật đầu: ``Quả đúng như thế nha-nanodesu!''

Sói đen-chan thì tỏ ra lo lắng với Thỏ-chan, vừa ngước nhìn pháo hoa vừa trông chờ:
``Mà\dots\ Pekora-chan sẽ ổn chứ?''

Tôi chỉ có thể tặc lưỡi, khoanh tay đáp: ``Chịu. Kệ em ấy thôi. Tự dưng lại chơi ngu làm
gì.'' Fubu-chan, kẻ đầu têu cho màn pháo hoa độc nhất vô nhị này lại giữ vẻ điềm tĩnh:
``Thôi nào, mọi người cứ tận hưởng màn pháo hoa siêu cấp vũ trụ này đi\~{}''

Koro-chan nhìn tiếp một quả pháo khổng lồ nổ bùng trên cao, liền cười khúc khích: ``Tôi
cũng muốn bay lên như Pekora-chan á\~{}'' Okayu cũng gật gù phụ họa: ``Tôi cũng thế\~{}''

Tôi chỉ biết thở dài với ý tưởng khá ngây ngô từ cặp đôi ``Cờ-hó và Mồn lèo'' này: ``Hai
đứa đúng là điên rồ thật.''

\dots Phải thừa nhận, màn trình diễn pháo hoa này trông còn hoành tráng và rực rỡ hơn so
với hồi ban nãy. Những màn pháo hoa bắn to hơn, nổ cũng vang hơn, lại còn nhiều màu
sắc hơn.

Tôi hy vọng là nó sẽ không gây sự chú ý cho toàn thành phố\dots\ Mà thôi kệ, có khi ai đó
đã quay lại màn trình diễn này rồi đăng lên mạng rồi cũng nên.

\ct

Màn pháo hoa kéo dài tận nửa tiếng, kết thúc bằng một hình ảnh không ai ngờ tới: khuôn
mặt của Pekora hiện ra giữa bầu trời, đang nháy mắt và giơ ngón tay cái đầy tự tin như
muốn nói: ``Tôi sẽ ổn thôi, peko!''

Nhưng chưa kịp thở phào, thân hình nhỏ bé của nàng Thỏ rơi thẳng xuống bệ phóng.

\textbf{*RẦM!*}

``Ui da! Peko!'' tiếng hét kêu đau của em ấy vang lên, vô tình kích hoạt cả đống pháo hoa
chưa kịp bắn. Chỉ chờ có thế, hàng loạt pháo lao đi mọi hướng: cái bay lên trời, cái bắn
xuống đất, cái thì thậm chí len lỏi vào các ô nhà, cái thì nhắm thẳng vào chúng tôi.

``\textbf{Chạy mau!!}'' Tôi hét lớn khi pháo hoa bắt đầu truy đuổi cả đám. Chúng tôi quyết định
chia tản ra để không gây thiệt hại nặng nhất.

Mà khổ nỗi, tôi lại chạy cùng với hai cô gái đại ngốc nhất của cả công ty này mới đau
chứ. Số tôi năm tới không biết liệu có đen không nữa\dots

\begin{dialogue}
    \speak{Roboco} Đừng có nhắm vào tôi chứ!!
    \speak{Noel} Lỗi kỹ thuật rồi! Lỗi kỹ thuật rồi!
    \speak{Kuon} Cái gì chứ!? Hai đứa sợ quá đổi hồn cho nhau luôn rồi à!?
\end{dialogue}

Trong khi tất cả đang hoảng loạn bỏ chạy, tôi nhìn thấy Cáo trắng-chan lại đứng thản nhiên
giữa khung cảnh hỗn loạn, nhìn vào máy quay mà em đã chuẩn bị sẵn, mỉm cười ``Và đó
là cách chúng tôi đón Tết Nguyên đán với một bữa tiệc pháo hoa nhé mọi người\~{}''

Mio-chan và tôi lập tức kéo em ấy lại, rồi cùng với tất cả các thành viên nhảy xuống thang
tầng sáu (bằng cách nào đó đã được mở.)

Pháo hoa tiếp tục bắn như tên lửa, gây náo loạn khắp khu phố. A-san, đứng từ tầng năm,
chỉ biết thở dài với trò quái quỷ đầu năm mới của các nàng ``hề'': ``Như thế này là giết
người ta còn gì\dots''

Hậu quả là cả công ty và gia đình tôi phải đi đền bù thiệt hại cho cả khu phố. Đáng ra là
như thế, nhưng các nàng thần tượng của chúng tôi lại nghĩ ra chiêu bài quen thuộc - tia
lãng quên - để cứu nguy.

Ít nhất thì chúng tôi cũng sẽ không được lên báo vì lý do gây phá hoại, mà chắc là cũng
vì màn trình diễn pháo hoa bí ẩn đó\dots\ Tôi cho là vậy.

Gia đình tôi, chứng kiến toàn bộ sự việc, chỉ biết ngồi thở dài trong sốc nặng: ``Quả đúng
là một món quà đầu năm thật `ý nghĩa' và `trân trọng' đến từ các gái nhà \hll!''


\end{document}